\documentclass[11pt,reqno]{amsart}
\usepackage[localtheorems]{raeez-math-template}

% ================================================================
% Packages
% ================================================================
\usepackage{array}

% ================================================================
% Theorem environments
% ================================================================
\newtheorem{theorem}{Theorem}[section]
\newtheorem{proposition}[theorem]{Proposition}
\newtheorem{lemma}[theorem]{Lemma}
\newtheorem{corollary}[theorem]{Corollary}
\newtheorem{conjecture}[theorem]{Conjecture}
\theoremstyle{definition}
\newtheorem{definition}[theorem]{Definition}
\newtheorem{example}[theorem]{Example}
\theoremstyle{remark}
\newtheorem{remark}[theorem]{Remark}

% ================================================================
% Macros
% ================================================================
\newcommand{\cA}{\mathcal{A}}
\newcommand{\cB}{\mathcal{B}}
\newcommand{\cC}{\mathcal{C}}
\newcommand{\cD}{\mathcal{D}}
\newcommand{\cF}{\mathcal{F}}
\newcommand{\cM}{\mathcal{M}}
\newcommand{\cN}{\mathcal{N}}
\newcommand{\cO}{\mathcal{O}}
\newcommand{\cP}{\mathcal{P}}
\newcommand{\cW}{\mathcal{W}}
\newcommand{\cZ}{\mathcal{Z}}
\newcommand{\barB}{\bar{B}}
\newcommand{\barBch}{\bar{B}^{\mathrm{ch}}}
\newcommand{\Ran}{\mathrm{Ran}}
\newcommand{\MC}{\mathrm{MC}}
\newcommand{\Sym}{\mathrm{Sym}}
\newcommand{\Symch}{\mathrm{Sym}^{\mathrm{ch}}}
\newcommand{\Hom}{\mathrm{Hom}}
\newcommand{\RHom}{R\!\operatorname{Hom}}
\newcommand{\Ext}{\mathrm{Ext}}
\newcommand{\Res}{\mathrm{Res}}
\newcommand{\FM}{\overline{C}}
\newcommand{\fg}{\mathfrak{g}}
\newcommand{\fh}{\mathfrak{h}}
\newcommand{\Vir}{\mathrm{Vir}}
\newcommand{\KM}{\mathrm{KM}}
\newcommand{\Heis}{\mathrm{Heis}}
\newcommand{\ChirHoch}{\mathrm{ChirHoch}}
\newcommand{\bC}{\mathbb{C}}
\newcommand{\bZ}{\mathbb{Z}}
\newcommand{\bR}{\mathbb{R}}
\newcommand{\bP}{\mathbb{P}}
\newcommand{\Einf}{E_\infty}
\newcommand{\Etwo}{E_2}
\newcommand{\Eone}{E_1}
\newcommand{\SCchtop}{\mathrm{SC}^{\mathrm{ch,top}}}
\newcommand{\Ainfty}{A_\infty}
\newcommand{\Linfty}{L_\infty}

\DeclareMathOperator{\gr}{gr}
\DeclareMathOperator{\av}{av}
\DeclareMathOperator{\coLie}{coLie}
\DeclareMathOperator{\Conf}{Conf}

\numberwithin{equation}{section}

% ================================================================
\begin{document}

\title[Koszulness chart functors and conditional atlas criteria]
{Koszulness chart functors and conditional atlas criteria:
the Kontsevich--Tamarkin reformulation of chiral Koszul duality}

\author{Raeez Lorgat}
\address{Perimeter Institute for Theoretical Physics,
 Waterloo, ON N2L 2Y5, Canada}
\email{rlorgat@perimeterinstitute.ca}

\date{\today}

\begin{abstract}
For a finite-PBW chiral algebra $\cA$ on an algebraic curve~$X$, the
Kontsevich--Tamarkin action organizes Koszulness tests into
associator-dependent chart functors. After choosing a framing and a
finite PBW-Artin stage, the Koszul locus is represented by an open
subfunctor of the finite-stage parameter space; on nonempty framed
components the chart transitions form the usual pro-\(\mathrm{GRT}_1\)
transfer torsor. This paper records the chart formulation and keeps
the logical directions separate. The algebraic/PBW core gives the
bidirectional Koszulness criteria. Hochschild, factorization-homology,
Lagrangian, and mixed-Hodge criteria are consequences or conditional
converses under the hypotheses stated in the theorem. The scalar
degree-two saturation statement is proved only on the scalar named
rows; outside those rows it remains a conditional chart assertion.
\end{abstract}

\maketitle

\tableofcontents

% ================================================================
\section{Introduction: the Koszul locus and why it matters}
\label{sec:introduction}
% ================================================================

\subsection{The problem}

The bar complex of a chiral algebra $\cA$ on an algebraic curve~$X$
is the cofree $\Eone$-coalgebra
\begin{equation}\label{eq:bar-complex}
\barBch(\cA) \;=\; T^c(s^{-1}\bar\cA)
\end{equation}
with deconcatenation coproduct, where $\bar\cA = \ker(\varepsilon)$
is the augmentation ideal and
$|s^{-1}v| = |v| - 1$ is the desuspension grading.
The bar differential extracts OPE data through the logarithmic
propagator $\eta = d\log(z_1 - z_2)$ on Fulton--MacPherson
configuration spaces $\FM_n(X)$.
This construction exists for every chiral algebra. Its differential
squares to zero. Its cohomology defines a factorization coalgebra.
The cobar functor $\Omega^{\mathrm{ch}}$ produces from
$\barBch(\cA)$ a new chiral algebra, and the counit
$\varepsilon_\cA \colon \Omega^{\mathrm{ch}}(\barBch(\cA)) \to \cA$
is a natural transformation.

The question is: when is $\varepsilon_\cA$ a quasi-isomorphism?

When it is, the bar complex is a complete invariant: it encodes~$\cA$
up to quasi-isomorphism, the spectral sequences computing its
cohomology collapse, the resolution is linear in the PBW degree,
and a single scalar $\kappa(\cA)$ controls the universal
Maurer--Cartan element at degree two. When it is not, the
bar--cobar object persists only in the coderived category, the
spectral sequences have nontrivial higher differentials, and the
resolution is nonlinear in a way that resists finite description.

The locus on which $\varepsilon_\cA$ is a quasi-isomorphism is the
\emph{Koszul locus}. This paper gives fourteen tests for this locus.
The algebraic/PBW tests are bidirectional criteria. The Hochschild,
factorization-homology, Lagrangian, and mixed-Hodge tests enter with
the hypotheses and logical directions stated below.

\subsection{Context}

In classical algebra, the Koszul property of a quadratic algebra
$A = T(V)/(R)$ has been characterized by
Priddy~\cite{Priddy} (purity of the bar complex),
Beilinson--Ginzburg--Soergel~\cite{BGS} (Ext diagonal vanishing),
and Loday--Vallette~\cite{LV12} ($\Ainfty$-formality of the
Koszul dual cooperad). Each formulation detects the same
phenomenon from a different mathematical vantage point: the
bar complex has diagonal bar cohomology exactly when it is
concentrated in degree one.

Chiral algebras are not quadratic. The Virasoro algebra has a
quartic OPE pole; $\cW$-algebras have poles of
arbitrarily high order; Yangians have spectral-parameter
relations. None of these admit a quadratic presentation, and the
Beilinson--Ginzburg--Soergel theory does not see them directly.
Three features of chiral algebras on algebraic curves defeat the
classical framework simultaneously: (a)~the bar differential on
the genus expansion satisfies
$d^2_{\mathrm{fib}} = \kappa(\cA) \cdot \omega_g$, so the
classical bar--cobar adjunction no longer applies beyond genus
zero; (b)~the tensor power $V^{\otimes n}$ is replaced
by the factorization product $\cA^{\boxtimes n}$ over $\FM_n(X)$,
and the differential is a sum of OPE residues at boundary strata,
not a combinatorial sum of face maps; (c)~conformal-weight
components are infinite-dimensional, so convergence of spectral
sequences is a separate problem.

The structure of chiral objects forces the extension of Koszul
theory beyond the classical quadratic setting. The
Francis--Gaitsgory adjunction
$\mathrm{Lie}^{\mathrm{ch}} \dashv \mathrm{Com}^{\mathrm{ch}}$
on $\Ran(X)$~\cite{FG12} settles Koszul duality of the
Lie--Com pair at the operadic level, but it does not determine
whether a specific chiral algebra is Koszul: it does not see the
genus corrections in the bar complex, and it does not detect the
shadow obstruction tower whose depth varies from two to infinity
across the standard families.

\subsection{The fourteen characterizations}

The main theorem of this paper assembles fourteen tests for
chiral Koszulness. They sort into four groups by mathematical
origin.

\smallskip
\noindent
\textbf{Homological algebra} (from the bar complex and its
spectral sequences): chiral Koszulness~(i), PBW
$\Etwo$-collapse~(ii), $\Ainfty$-formality~(iii), Ext
diagonal vanishing~(xii), Kac--Shapovalov determinant
nonvanishing~(xiii).

\smallskip
\noindent
\textbf{Deformation theory} (from the convolution algebra and
its formality): shadow-tower formality~(iv),
$\Etwo$-formality of chiral Hochschild cohomology~(v),
Sklyanin $H^2 = 0$~(xiv).

\smallskip
\noindent
\textbf{Factorization geometry} (from configuration spaces and
the Ran space): curve independence~(vi), PBW
universality~(vii), Barr--Beck--Lurie monadicity~(viii),
factorization homology concentration~(ix),
FM boundary acyclicity~(x).

\smallskip
\noindent
\textbf{Tropical geometry} (from the dual complex of the FM
compactification): tropical Koszulness~(xi).

\smallskip
\noindent
Two further conditions stand outside these fourteen. The
\emph{Lagrangian criterion}~(xv): $\cM_\cA$ and $\cM_{\cA^!}$
are transverse Lagrangians in the $(-1)$-shifted symplectic
deformation space, proved equivalent under a perfectness
hypothesis. And \emph{$\cD$-module purity}~(xvi): each
$\barBch_n(\cA)$ is pure of weight~$n$ as a mixed Hodge module,
proved to imply Koszulness with the converse established for
the affine Kac--Moody family and open in general.

\subsection{The six-fold TFAE on the Koszul locus, with a class-$G$
seventh equivalence: hub-and-spoke reconstitution}
\label{subsec:seven-fold-tfae-hub-spoke}

The fourteen characterizations above form an atlas of
Koszulness chart functors. At genus zero, a strictly smaller
sub-atlas admits a sharper presentation: six characterizations
are mutually equivalent through a single hub on the full Koszul
locus, and a seventh (SC-formality) is equivalent to the same hub
when attention is restricted to the class-$G$ stratum.

\begin{theorem}[Six-fold TFAE at genus zero, with a class-$G$
seventh equivalence; hub-and-spoke]\label{thm:seven-fold-tfae-hub-spoke-sa}
Let $\cA$ be an augmented chiral algebra on a smooth genus-zero
curve $X$ equipped with a PBW filtration $F_\bullet$, and let
$\tau_{\mathrm{univ}}\colon \barBch(\cA) \to \cA$ be the
universal twisting morphism of
the canonical bar twisting morphism
$\tau_{\mathrm{univ}}\colon\barBch(\cA)\to\cA$. Conditions
\textup{\ref{sv1}--\ref{sv6}} below are equivalent \emph{through
the PBW $\Etwo$-collapse hub} on the full Koszul locus;
condition \textup{\ref{sv7}} is equivalent to \textup{\ref{sv6}}
on the class-$G$ stratum and otherwise fails.
\begin{enumerate}[label=\textup{(S\arabic*)},leftmargin=3em]
\item\label{sv1}
  \emph{Chiral Koszul morphism}: $\tau_{\mathrm{univ}}$ is a
  chiral Koszul morphism (Definition
  \textup{\ref{def:kfc-chiral-koszul-morphism}}).
\item\label{sv2}
  \emph{Counit quasi-isomorphism}: the counit
  $\varepsilon_{\tau}\colon
   \Omega_X(\barBch(\cA)) \to \cA$ is a quasi-isomorphism.
\item\label{sv3}
  \emph{Unit weak equivalence}: the unit
  $\eta_\tau\colon \barBch(\cA) \to
   \barBch(\Omega_X(\barBch(\cA)))$ is a weak equivalence of
  conilpotent complete factorisation coalgebras.
\item\label{sv4}
  \emph{Twisted tensor acyclicity}: both twisted tensor
  products $K_\tau^L(\cA, \barBch(\cA))$ and
  $K_\tau^R(\barBch(\cA), \cA)$ are acyclic.
\item\label{sv5}
  \emph{Bar concentration in weight~$1$}: bar cohomology
  $H^{p,q}(\barBch(\cA))$ vanishes for $q \neq 0$, with
  $H^{p,0}$ concentrated in bar-weight $p$.
\item\label{sv6}
  \emph{PBW $\Etwo$-collapse} \textup{(the hub)}: the PBW
  spectral sequence of $\barBch(\cA)$ collapses at $E_2$:
  $d_r = 0$ for all $r \geq 2$.
\item\label{sv7}
  \emph{\textup{SC}-formality at genus zero}
  \textup{(class-G scoped)}: the transferred Swiss-cheese
  operation tower has $m_k^{\mathrm{SC}} = 0$ for all
  $k \geq 3$.
\end{enumerate}

\textup{(Hub-and-spoke structure.)} The hub is \textup{\ref{sv6}}.
For each $i \in \{1,2,3,4,5\}$ there is a single proved
bidirection $i \Leftrightarrow \textup{\ref{sv6}}$; the six-fold
TFAE on the Koszul locus follows by transitivity from these five
spokes. A sixth, class-$G$-scoped spoke
$\textup{\ref{sv7}} \Leftrightarrow \textup{\ref{sv6}}$ furnishes
the seventh equivalence on that stratum.

\textup{(Load-bearing bidirection.)} The spoke
$\textup{\ref{sv4}} \Leftrightarrow \textup{\ref{sv6}}$
(twisted-tensor $\Leftrightarrow$ PBW) is the substantive
bidirection: both forward and reverse require
genuine theorems (PBW implies twisted-tensor acyclicity through
the free resolution; twisted-tensor acyclicity implies PBW
through degeneration of the filtration spectral sequence). The
affine Kac--Moody family $V_k(\mathfrak{sl}_2)$ at generic $k$
witnesses both implications concretely.

\textup{(Class-G restriction of \textup{\ref{sv7}}.)} Spoke
\textup{\ref{sv7}} is class-G scoped: it is equivalent to
\textup{\ref{sv6}} only when~$\cA$ is class~G. Class~L, C, M
algebras are chirally Koszul \textup{(}all six
non-\textup{\ref{sv7}} spokes hold\textup{)} yet fail
\textup{SC}-formality \textup{(}compare
Remark~\textup{\ref{warn:koszulness-not-sc-formality}})}.

\textup{(Higher-genus extension.)} Spokes
\textup{\ref{sv1}--\ref{sv4}} survive verbatim to $g \geq 1$
within the Koszul locus; spoke \textup{\ref{sv5}} survives to
$g \geq 1$ only for class~G, since
$d^2_{\mathrm{fib}} = \kappa \cdot \omega_g$ obstructs
weight-$1$ concentration off class~G.
\end{theorem}

\begin{proof}[Proof sketch]
The six spokes are proved elsewhere: spoke \textup{\ref{sv1}} by
Definition~\textup{\ref{def:kfc-chiral-koszul-morphism}} together
with the Koszul characterisation through PBW $\Etwo$-collapse;
spoke \textup{\ref{sv2}} is the bar--cobar counit criterion,
which is equivalent to $\Etwo$-collapse through the twisting
morphism dictionary; spoke \textup{\ref{sv3}} by the
conilpotent-complete universal property of
$\barBch$; spoke \textup{\ref{sv4}} by the substantive PBW
bidirection cited above; spoke \textup{\ref{sv5}} by
Priddy's criterion in its chiral avatar; spoke \textup{\ref{sv7}}
by the class-$G$ SC-formality criterion in the parent monograph.
The transitivity closure is the elementary observation
that six proved bidirections through a common hub determine the
remaining $15$ bidirections of the $21$-arrow complete graph.
\end{proof}

\begin{remark}[Why the hub is PBW $\Etwo$-collapse]
\label{rem:hub-is-pbw-e2}
The PBW spectral sequence is the universal invariant that
compares the associated graded algebra of~$\cA$ (determined by
the filtration alone, independent of the OPE) with the bar
cohomology of~$\cA$ (determined by the OPE, with the filtration
tracking PBW depth). Its $E_2$-collapse is the locus where these
two pieces of data agree. Every classical criterion reduces to
this comparison; all sit at a distance of one spoke from the hub.
\end{remark}

\subsection{Two distinctions}

\begin{remark}[Shadow depth and Koszulness]
\label{warn:shadow-depth-not-koszulness}
The four complexity classes G/L/C/M, with shadow depths
$\{2, 3, 4, \infty\}$, record the degree at which the shadow
obstruction tower first becomes nontrivial. They do \emph{not}
record failure of Koszulness. Every standard chiral algebra
is chirally Koszul. The Virasoro algebra has shadow depth
$r_{\max} = \infty$ (class~M); it is chirally Koszul at
generic central charge. Shadow depth classifies complexity
\emph{within} the Koszul world.
\end{remark}

\begin{remark}[Koszulness and SC-formality]
\label{warn:koszulness-not-sc-formality}
Koszulness is the condition that bar cohomology is concentrated
in degree one; SC-formality is the condition that the
higher operations $m_k^{\mathrm{SC}} = 0$ for $k \geq 3$ in
the Swiss-cheese bar complex. All standard families are chirally
Koszul; only class~G (Heisenberg) is SC-formal.
Confusing the two misidentifies what the shadow tower measures.
\end{remark}

\subsection{Conventions}

All chain complexes are cohomologically graded: $|d| = +1$.
Desuspension lowers degree: $|s^{-1}v| = |v| - 1$. The bar
complex uses the augmentation ideal $\bar\cA = \ker(\varepsilon)$;
we write $\barBch(\cA) = T^c(s^{-1}\bar\cA)$ with deconcatenation
coproduct. We work in characteristic zero over an algebraically
closed field. The Fulton--MacPherson compactification $\FM_n(X)$
is the real oriented blowup along all diagonals; it is a manifold
with corners whose boundary strata are indexed by rooted
trees~\cite{FultonMacPherson94}. We cite~\cite{LorgatMKD1} for background
results established in the parent monograph.


% ================================================================
\section{Koszulness chart functors}
\label{sec:moduli}
% ================================================================

\subsection{Principle}

The classical literature presents Koszulness through lists of
equivalent conditions. That presentation identifies a property with
the charts that detect it. In the chiral setting each chart is tied
to a choice of Drinfeld associator. Tamarkin transfer carries a
chart at one associator to a chart at another. The natural object is
therefore a chart functor: finite-stage algebraic on PBW-Artin
truncations, and pro-algebraic only after passing to the full
associator torsor.

\begin{definition}[Koszulness chart functor]
\label{def:koszulness-moduli-scheme}
Let $\cA$ be a chiral algebra on~$X$ with PBW filtration, and fix
a finite PBW-Artin stage~$N$ together with a framing of its
associated graded. The \emph{Koszulness chart functor}
$\cM_{\mathrm{Kosz}}^{\leq N,\mathrm{fr}}(\cA)$ is the functor
\[
\cM_{\mathrm{Kosz}}^{\leq N,\mathrm{fr}}(\cA)\colon
\mathbf{CAlg}_\bQ \longrightarrow \mathbf{Set},
\qquad
R \longmapsto
\bigl\{(\Phi, c_\Phi) :
\Phi \in \mathrm{Assoc}(R),\;
c_\Phi \in \mathrm{Chart}_\Phi(\cA;R)\bigr\}
\]
where $\mathrm{Assoc}(R)$ is the torsor of Drinfeld
associators~\cite{Drinfeld90} over $R$, and
$\mathrm{Chart}_\Phi(\cA;R)$ is the set of \emph{chart coordinates}:
triples $(\Pi_\Phi, \mathrm{wit}_\Phi, \gamma_\Phi)$ consisting of
\begin{enumerate}[label=\textup{(\alph*)}]
\item a Koszulness predicate
  $\Pi_\Phi$ depending only on the pair
  $(\cA, \Phi)$;
\item a witness $\mathrm{wit}_\Phi(\cA)$ in an appropriate test
  category, providing a proof witness for
  $\Pi_\Phi(\cA)$;
\item a transition 1-cocycle
  $\gamma_\Phi \colon \Phi \leadsto \Phi'$
  to every other associator, satisfying the Tamarkin
  compatibility~\cite{Tamarkin03}
  \[
  \Pi_\Phi(\cA) \iff \Pi_{\Phi \cdot g}(\cA)
  \qquad \forall\,g \in \mathrm{GRT}_1(R).
  \]
\end{enumerate}
\end{definition}

\begin{theorem}[Finite-stage Koszulness atlas]
\label{thm:koszulness-moduli}
\textup{(Kontsevich--Tamarkin reformulation.)}
Let $\cA$ be a positive-energy chiral algebra on~$X$ with finite
PBW windows. After choosing a finite PBW-Artin stage and a framing,
the Koszul locus is represented by an open subfunctor of the
finite-stage parameter space. On every nonempty framed component,
the transition functors between associator charts form the usual
pro-\(\mathrm{GRT}_1\) transfer torsor. More precisely:
\begin{enumerate}[label=\textup{(A\arabic*)},leftmargin=3em]
\item\label{item:moduli-A1}
  The finite-stage chart functor is nonempty precisely when the
  chosen stage contains a witness for the chiral Koszul morphism of
  Definition~\textup{\ref{def:kfc-chiral-koszul-morphism}}.
  Passing to the inverse limit over PBW-Artin stages recovers
  chiral Koszulness under the complete separated
  Mittag--Leffler hypothesis.
\item\label{item:moduli-A2}
  The forgetful map
  $\pi\colon
  \cM_{\mathrm{Kosz}}^{\leq N,\mathrm{fr}}(\cA)
  \to \mathrm{Assoc}$
  carries a pro-\(\mathrm{GRT}_1\)-transfer action on each
  nonempty framed component.
\item\label{item:moduli-A3}
  The fourteen classical characterizations
  \textup{(i)}--\textup{(xiv)} give a finite chart atlas
  $\{U_j\}_{j=1}^{14}$: charts $U_1,\ldots,U_{10}$ are
  centred at the KZ associator $\Phi_{\mathrm{KZ}}$;
  $U_{11}$ is centred at $\Phi_{\mathrm{AT}}$;
  $U_{12}$ at $\Phi_{\mathrm{dR,B}}$;
  $U_{13}$ at $\Phi_{\mathrm{ell}}$;
  $U_{14}$ at $\Phi_{\mathrm{Kon}}$.
  Intersections are identified by Tamarkin transfer whenever both
  charts contain a Koszul witness.
\item\label{item:moduli-A4}
  The shadow-class stratification
  $G \sqcup L \sqcup C \sqcup M \sqcup FF$ \textup{(}Feigin--Frenkel
  critical stratum\textup{)} selects a distinguished chart on each
  stratum: $\Phi_{\mathrm{KZ}}$ on~$G$ and on the generic
  Kac--Moody locus within~$L$, $\Phi_{\mathrm{ell}}$ on~$L$,
  $\Phi_{\mathrm{AT}}$ on~$C$, $\Phi_{\mathrm{Kon}} = \Phi_\infty$
  on~$M$, and $\Phi_{\mathrm{dR,B}}$ on~$FF$.
\end{enumerate}
\end{theorem}

\begin{proof}
\emph{\ref{item:moduli-A1}.}
Definition~\ref{def:kfc-chiral-koszul-morphism} is
$\mathrm{GRT}_1$-invariant: acyclicity of the twisted tensor
product $K_\tau^L(\cA,\barBch(\cA))$ is preserved under the
Tamarkin transfer $K_\tau^L \mapsto K_{\Phi \cdot \tau}^L$,
because the transfer is a quasi-isomorphism
(Kontsevich~\cite{Kontsevich99}, Tamarkin~\cite{Tamarkin03}).
At a fixed finite PBW-Artin stage this gives an open condition on
the framed parameter space: the determinant minors for the two
twisted tensor complexes must be invertible. Taking the inverse
limit over finite windows gives the full Koszul predicate when the
PBW filtration is complete, separated, and satisfies the
Mittag--Leffler condition.

\emph{\ref{item:moduli-A2}.}
Fix $\Phi_0 = \Phi_{\mathrm{KZ}}$ as base point. For any other
$\Phi$, the Grothendieck--Teichm\"uller torsor provides a unique
$g_{\Phi} \in \mathrm{GRT}_1$ with $\Phi = \Phi_0 \cdot g_\Phi$.
The Tamarkin transfer
$g_\Phi^* \colon \mathrm{Chart}_{\Phi_0}(\cA) \to
\mathrm{Chart}_\Phi(\cA)$
is an equivalence on each framed component containing a Koszul
witness, and the cocycle condition
$g_{\Phi''} = g_{\Phi'} \cdot h_{\Phi' \to \Phi''}$
holds because associator composition is strictly associative in
$\mathrm{GRT}_1$. Thus the chart transitions carry the expected
pro-\(\mathrm{GRT}_1\) torsor structure. The global object is a
pro-algebraic chart functor assembled from finite PBW stages.

\emph{\ref{item:moduli-A3}.}
The ten classical characterizations of
Theorem~\ref{thm:kfc-fourteen}~(i)--(x) sit at $\Phi_{\mathrm{KZ}}$:
each is stated in the KZ language (PBW spectral sequence, bar
cohomology, $\Ainfty$-formality, BBL monadicity, etc.) whose
coherence data is the KZ associator. The four supplementary
charts $U_{11},\ldots,U_{14}$ are constructed in
Propositions~\ref{prop:at-chart}, \ref{prop:hodge-betti-chart},
\ref{prop:elliptic-chart}, and~\ref{prop:kontsevich-chart}
below, with the hypotheses stated in those propositions.
Non-emptiness of a pairwise intersection means that both charts
contain a Koszul witness for the same finite PBW stage.
If $\Phi_2=\Phi_1\cdot g$, Tamarkin transfer by
$g\in\mathrm{GRT}_1$ identifies the two chart functors on that
component.

\emph{\ref{item:moduli-A4}.}
The shadow class is an intrinsic invariant of the MC element
$\Theta_\cA$; the selection of a distinguished chart on each
stratum is determined by the computational complexity of the
shadow algebra bracket $[-,-]^{\mathrm{sh}}$ there.
Class~G: abelian shadow algebra, KZ associator is the
canonical flat connection on $\FM_n(\bC)$, and~$U_1$ (PBW
$\Etwo$-collapse) is trivially satisfied.
Class~L: cubic bracket only, genus-one Enriquez coupling; the
elliptic associator $\Phi_{\mathrm{ell}}$ with its modular
equivariance is the natural chart.
Class~C: contact termination at $S_4$, AT Lagrangian coordinates
diagonalize the quartic bracket (Proposition~\ref{prop:at-chart}).
Class~M: infinite tower, Kontsevich integral gives chart
coordinates tracking the full MC expansion
(Proposition~\ref{prop:kontsevich-chart}).
Class~FF: Feigin--Frenkel locus $k = -h^\vee$; the de~Rham--Betti
chart $\Phi_{\mathrm{dR,B}}$ sees the purity via BGG resolution
(Proposition~\ref{prop:hodge-betti-chart}).
\end{proof}

\begin{remark}[Kontsevich--Tamarkin reformulated]
\label{rem:kt-reformulation}
The classical Kontsevich--Tamarkin theorem~\cite{Kontsevich99,Tamarkin03}
is usually stated: ``formality of little disks holds universally,
but the choice of formality isomorphism depends on the associator.''
The Koszulness chart functor gives this statement operational
content at the level of a single algebra: chiral Koszulness is a
$\mathrm{GRT}_1$-invariant property of $\cA$; the classical
characterizations are $\mathrm{GRT}_1$-coordinate charts on the
finite PBW-Artin chart functors.
Switching between two classical criteria is an application of the
Tamarkin $\Phi$-transfer; the transfer becomes trivial on the
$\mathrm{GRT}_1$-invariants.
\end{remark}

\subsection{Atlas}

The atlas has fourteen charts. The first ten are the classical
conditions (the old Theorem 2.2, stated below with its original
numbering as Theorem~\ref{thm:kfc-fourteen}); the additional four
are finite-PBW comparison charts. They are independent chart
coordinates, but their converses require the perfectness, purity,
genus-one flatness, or two-colour operadic hypotheses stated below.

\begin{proposition}[$U_{11}$: Alekseev--Torossian chart]
\label{prop:at-chart}
At the Alekseev--Torossian associator
$\Phi_{\mathrm{AT}}$~\cite{AT12}, chiral Koszulness of a
chiral algebra~$\cA$ with PBW filtration is equivalent to
Lagrangian transversality of $\cM_\cA$ and $\cM_{\cA^!}$ in the
$(-1)$-shifted symplectic AT deformation space
$\cM^{\mathrm{AT}}_{\mathrm{comp}}$, provided the ambient tangent
complex is two-term perfect and the AT transfer preserves the
finite PBW filtration.
\end{proposition}

\begin{proof}[Proof sketch]
The Alekseev--Torossian associator is constructed from the
Kashiwara--Vergne equations on free Lie algebras~\cite{AT12}: it
provides an explicit homotopy between the KZ associator and its
symmetrization, together with a canonical choice of
Baker--Campbell--Hausdorff basis on the target.
Transfer the $(-1)$-shifted symplectic deformation space
$\cM_{\mathrm{comp}}$ of the classical Lagrangian criterion
to $\cM^{\mathrm{AT}}_{\mathrm{comp}}$ using the AT gauge:
\[
\Psi_{\mathrm{AT}} \colon
\cM_{\mathrm{comp}} \xrightarrow{\;\sim\;}
\cM^{\mathrm{AT}}_{\mathrm{comp}},
\qquad
\Psi_{\mathrm{AT}}^*\omega_{\mathrm{AT}} = \omega,
\]
where $\omega, \omega_{\mathrm{AT}}$ are the respective
symplectic forms. The AT gauge is constructed so that the ambient
tangent complex of $\cM^{\mathrm{AT}}_{\mathrm{comp}}$ is
two-term perfect: concentrated in degrees $\{1,2\}$, with
generators the Kashiwara--Vergne $F, H$ pair and relations the
Duflo--Kontsevich isomorphism. The chart is therefore available on
the finite-PBW component where this two-term perfect model exists.
Lagrangian transversality of $\Psi_{\mathrm{AT}}(\cM_\cA)$ and
$\Psi_{\mathrm{AT}}(\cM_{\cA^!})$ is equivalent to acyclicity of
the AT-twisted tensor product
$K_\tau^{\mathrm{AT}}(\cA, \cA^!)$, which equals the KZ-twisted
product up to a quasi-isomorphism induced by the AT gauge; hence
Koszulness.
\end{proof}

\begin{proposition}[$U_{12}$: Hodge--Betti chart]
\label{prop:hodge-betti-chart}
At the de~Rham--Betti associator
$\Phi_{\mathrm{dR,B}}$~\cite{Deligne89}, chiral Koszulness of a
chiral algebra~$\cA$ with conformal weight grading is equivalent
to $\cD$-module purity of each $\barBch_n(\cA)$ as a mixed Hodge
module~\textup{(}condition~\textup{(xvi)}\textup{)} in the
families where the BGG or DS comparison supplies a pure finite
PBW resolution. In general, purity implies FM boundary acyclicity;
the converse remains a family-specific assertion.
\end{proposition}

\begin{proof}[Proof sketch]
The de~Rham--Betti associator sits at the boundary point of the
Drinfeld associator torsor where the KZ connection on
$\FM_n(\bC)$ is identified with the comparison isomorphism
between de~Rham and Betti cohomology of a mixed Hodge
structure~\cite{Deligne89}. At this base point, chart
coordinates are weight-filtered mixed Hodge structures on the
bar complex.

For any $\cA$ with polynomial PBW associated graded, the
Hodge--Betti chart detects purity directly: the weight filtration
on $\barBch_n(\cA)$ is recovered from the Hodge--Betti period
matrix, and purity means the weight-filtered piece agrees with
the pure part. The implication
$(\mathrm{xvi}) \Rightarrow (\mathrm{x})$ is classical
(Saito's theorem on mixed Hodge modules). The converse is proved
only where an independent pure resolution is available:
for $\cA = \Vir_c$ at generic~$c$, the BGG resolution of the
vacuum Verma module~\cite{FeiginFuchs90,RCW82}
provides a pure resolution of the bar complex; combined with
non-degeneracy $\det G_h \neq 0$ outside the discrete
minimal-model lattice, this gives purity directly. The same
argument applies to the simple $\cW$-rows for which the DS
reduction preserves the relevant pure PBW resolution.
\end{proof}

\begin{proposition}[$U_{13}$: Enriquez elliptic chart]
\label{prop:elliptic-chart}
At the Enriquez elliptic associator
$\Phi_{\mathrm{ell}}$~\cite{Enriquez14}, chiral Koszulness of
$\cA$ is equivalent to flatness of the genus-one twisted
K\"unneth formula
$\barBch(\cA \boxtimes_{\mathrm{ell}} \cA^!) \simeq
 \barBch(\cA) \boxtimes \barBch(\cA^!)$
on the elliptic configuration space $\overline{\cM}_{1,n}$.
This is a genus-one chart criterion on the finite PBW component
where the KZB transfer and the elliptic PBW spectral sequence are
defined.
\end{proposition}

\begin{proof}[Proof sketch]
The Enriquez elliptic associator~\cite{Enriquez14} is
constructed from the KZB connection on the universal elliptic
curve; it provides a canonical elliptic completion of the KZ
associator, with modular group $\mathrm{SL}_2(\bZ)$-equivariance
built in. The chart coordinate at $\Phi_{\mathrm{ell}}$ is the
elliptic twisted K\"unneth isomorphism: for chirally Koszul
$\cA$, the genus-one bar complex factorises as a K\"unneth
product of genus-zero data twisted by the elliptic R-matrix.

Flatness of the twisted K\"unneth is equivalent to
$\Etwo$-collapse of the genus-one PBW spectral sequence on the
elliptic curve. This is condition~(ii) transported to
$\FM_n(E)$; the classical proof on $\FM_n(\bC)$ carries over,
using the fact that the elliptic R-matrix (Felder 1994) is the
image of the rational R-matrix under the Tamarkin transfer
from $\Phi_{\mathrm{KZ}}$ to $\Phi_{\mathrm{ell}}$.
Hence $U_{13}$ is witnessed on the elliptic finite-PBW component;
the chart
distinguishes the class~L stratum because affine Kac--Moody
algebras at generic level are where the genus-one Enriquez
coupling first activates (cubic bracket~$S_3$ on the shadow
algebra).
\end{proof}

\begin{proposition}[$U_{14}$: Kontsevich integral chart]
\label{prop:kontsevich-chart}
At the Kontsevich integral associator
$\Phi_{\mathrm{Kon}}$~\cite{Kontsevich93}, chiral Koszulness of
$\cA$ is equivalent to homotopy-Koszulity of the Swiss-cheese
operad $\SCchtop$ acting on the pair
$(\cZ^{\mathrm{der}}_{\mathrm{ch}}(\cA), \cA)$ at every degree of
the MC expansion, assuming a two-colour chain-level
Kontsevich--Tamarkin formality comparison for the chosen
finite-PBW component.
On the class~M stratum, this chart records the infinite shadow
tower in chord-diagram coordinates.
\end{proposition}

\begin{proof}[Proof sketch]
The Kontsevich integral~\cite{Kontsevich93} is the universal
Vassiliev invariant: it assembles into an associator
$\Phi_{\mathrm{Kon}}$ at the Vassiliev base point, with chart
coordinates indexed by chord diagrams. For chiral Koszul~$\cA$,
the MC element $\Theta_\cA$ is a sum over chord diagrams in the
Kontsevich expansion, each coefficient computed by the shadow
algebra bracket.

Homotopy-Koszulity of $\SCchtop$ on the pair
$(\cZ^{\mathrm{der}}_{\mathrm{ch}}(\cA), \cA)$ is the condition
that the bar complex of $\SCchtop$ has cohomology concentrated
on the operadic diagonal. At the Kontsevich base point, this is
the assertion that chord diagram coefficients satisfy the
IHX and STU relations, which is automatic for~$\SCchtop$ by
Kontsevich formality for the little disks in two colors, once the
two-colour chain-level comparison is fixed.
On class~M, the Virasoro shadow tower is an infinite sum over chord
diagrams indexed by
the tower levels, and $\SCchtop$ homotopy-Koszulity is precisely
the statement that this sum converges in the coderived
category~\cite{PP}.
\end{proof}

\begin{remark}[Chart redundancy and the torsor structure]
\label{rem:chart-redundancy}
The fourteen charts are not equal in logical strength. The core
PBW charts give bidirectional criteria; the supplementary charts
give consequences or conditional converses. The content of
Theorem~\ref{thm:koszulness-moduli}~\ref{item:moduli-A2} is that
the overlap between valid charts is organized by the
$\mathrm{GRT}_1$-torsor structure, and the content of
\ref{item:moduli-A4} is that the redundancy is broken on each
stratum by the distinguished chart. Each chart is adapted to a
stratum, and finiteness here means finite after choosing the
named comparison base points, not finite presentation of the full
associator torsor.
\end{remark}

\subsection{Virasoro non-circular closure}

The classical proof of Virasoro Koszulness at generic~$c$ traced
a bar--cobar circle: Koszulness implies $\Etwo$-collapse implies
concentration of chiral Hochschild cohomology, and the converse
closes the loop through bar--cobar inversion. Every step uses the
bar complex; circularity is latent. The Hodge--Betti chart
$U_{12}$ provides an independent proof via Feigin--Fuchs BGG
resolution.

\begin{theorem}[Virasoro Koszulness, non-circular]
\label{thm:virasoro-koszulness-non-circular}
For the Virasoro vertex algebra $\Vir_c$ at generic central
charge $c \notin \{c_{p,q} : (p,q)\in\bZ^2_{>0},\gcd(p,q)=1\}$,
the bar cohomology weight components satisfy
\[
H^\bullet(\Vir_c, \Vir_c)_{\mathrm{weight}\,r+1} = 0
\qquad \text{for all } r \geq 2,
\]
using the Hodge--Betti comparison. Consequently,
$\Vir_c$ is chirally Koszul, proved directly at the $U_{12}$
chart.
\end{theorem}

\begin{proof}
We work entirely at the Hodge--Betti chart $U_{12}$. The proof
has two classical inputs.

\emph{Step 1: Feigin--Fuchs BGG resolution.}
For each highest weight~$h$ in the vacuum Verma module
$M_0(c)$ of $\Vir_c$, the Feigin--Fuchs theorem~\cite{FeiginFuchs90}
(see also Rocha-Caridi--Wallach~\cite{RCW82}) constructs a
resolution of the vacuum irreducible by Verma modules indexed by
elements of an affine Weyl group acting on weights:
\[
\cdots \to
\bigoplus_{w \in W^{(2)}} M_{w \cdot 0}(c) \to
\bigoplus_{w \in W^{(1)}} M_{w \cdot 0}(c) \to
M_0(c) \to L(c,0) \to 0,
\]
with differentials given by singular vector embeddings. At
generic~$c$ (outside the minimal-model lattice), each singular
vector embedding is nonzero and the resolution is exact. Each
term is a Verma module, whose mixed Hodge structure is pure of
weight given by its conformal weight (this is classical; Verma
modules are free over the PBW algebra).

\emph{Step 2: Kac determinant non-degeneracy.}
The Kac determinant~\cite{KacBook,FeiginFuchs90} at conformal
weight~$h$ is
\[
\det G_h
\;=\; \prod_{\substack{p, q \geq 1 \\ pq \leq h}}
(h - h_{p,q}(c))^{P(h - pq)},
\]
where $P$ is the partition function and
$h_{p,q}(c) = ((p\alpha_+ + q\alpha_-)^2 - (\alpha_+ + \alpha_-)^2)/4$
with $\alpha_\pm$ the Liouville parameters determined by~$c$.
For~$c$ outside the minimal-model lattice, the Kac determinant is
nonzero at every bar-relevant weight~$h$: the BGG resolution of
Step~1 embeds faithfully in the bar complex.

\emph{Step 3: Purity of the bar complex.}
Combine Steps~1 and~2. The BGG resolution of Step~1 gives a
$\cD$-module resolution of the unit~$\omega_X$ of $\Vir_c$ by a
sum of Verma modules; each Verma module is pure of its conformal
weight in the Hodge--Betti category (Verma modules have free PBW
associated graded, hence polynomial mixed Hodge structure). Step~2
(Kac determinant non-degeneracy at generic~$c$) guarantees that
the BGG differentials are injective at every weight, so the
resolution is exact. The bar complex $\barBch_n(\Vir_c)$ is
quasi-isomorphic to the $n$-fold iterate of this resolution,
and iterated pure objects remain pure by Saito's stability
theorem~\cite{Saito}.

\emph{Step 4: Weight vanishing.}
Purity of $\barBch_n(\Vir_c)$ of weight~$n$ forces
$H^\bullet(\Vir_c,\Vir_c)_{\mathrm{weight}\,r+1} = 0$ for every
$r \geq 2$ off the Koszul diagonal: any nonzero class in weight
$r + 1 \neq n$ would contradict purity of the bar complex in
weight~$n$ after the Ext computation.

\emph{Conclusion.} The four steps use only Feigin--Fuchs BGG
(Step~1), Kac determinant (Step~2), Saito stability (Step~3),
and purity vanishing (Step~4).
\end{proof}

\begin{remark}[Non-circular input]
\label{rem:fm197-closed}
The Hodge--Betti proof uses data external to bar--cobar inversion:
the Feigin--Fuchs BGG resolution and the Kac determinant. The bar
complex receives that input through purity of the induced
$\cD$-module resolution.
\end{remark}

\subsection{Yangian chart inclusion}

The associative Koszulness of the Yangian $Y_\hbar(\fg)$
uses the classical Polishchuk--Positselski PBW criterion. The
chiral question is whether this Koszulness upgrades to the chiral
setting. The chart-functor formulation isolates the precise
comparison map.

\begin{theorem}[Yangian chart inclusion]
\label{thm:yangian-chart-inclusion}
For the Yangian $Y_\hbar(\fg)$ with $\fg$ simple, assume that the
RTT evaluation data define a finite-PBW \(\Eone\)-chiral avatar and
that the chart-comparison functor is representable and reflects the
Shapovalov open locus. Then every such associative PBW chart maps to
the corresponding chiral Koszulness chart:
\[
\cM_{\mathrm{Kosz}}^{\mathrm{assoc}}(Y_\hbar)
\hookrightarrow
\cM_{\mathrm{Kosz}}^{\mathrm{chiral}}(Y_\hbar).
\]
This map is a $\mathrm{GRT}_1$-equivariant open immersion on the
comparison component. Chart-surjectivity remains open in higher
rank.
\end{theorem}

\begin{proof}[Proof of the embedding]
For $\Phi = \Phi_{\mathrm{KZ}}$, the associative PBW chart $U_1$
detects $Y_\hbar(\fg)$-Koszulness in the classical PBW sense.
Transport this chart coordinate to the chiral setting through the
assumed finite-PBW $\Eone$-chiral avatar: the RTT data defines the
required formal-disk chiral structure via the evaluation modules
$V_u$, and the PBW filtration is compatible with the chiral bracket.
The associative chart coordinate
$(\Phi_{\mathrm{KZ}}, c^{\mathrm{assoc}}_{\mathrm{KZ}})$ maps to
the chiral chart coordinate
$(\Phi_{\mathrm{KZ}}, c^{\mathrm{chiral}}_{\mathrm{KZ}})$: both
detect $\Etwo$-collapse of the same PBW spectral sequence, the
associative one on $Y_\hbar(\fg)$ and the chiral one on its
$\Eone$-chiral avatar. Hence
$\cM_{\mathrm{Kosz}}^{\mathrm{assoc}}(Y_\hbar) \hookrightarrow
 \cM_{\mathrm{Kosz}}^{\mathrm{chiral}}(Y_\hbar)$
is defined; $\mathrm{GRT}_1$-equivariance of both sides makes it
an equivariant map.

By hypothesis, the associative chart~$U_1^{\mathrm{assoc}}$ is the
complement of the zero locus of the Shapovalov determinant, and its
preimage in $\cM_{\mathrm{Kosz}}^{\mathrm{chiral}}(Y_\hbar)$ is the
complement of the same zero locus. Hence the comparison morphism is
open on this component.
\end{proof}

\begin{remark}[Associative and chiral Yangian charts]
\label{rem:fm161-closed}
The chart formulation separates associative Koszulness from
chiral Koszulness. The associative atlas embeds in the chiral
atlas on the comparison component. For
$\fg=\mathfrak{sl}_2$ the embedding is an isomorphism; for higher
rank, chart-surjectivity is the remaining open question.
\end{remark}

% ================================================================
\section{The fourteen characterizations}
\label{sec:fourteen}
% ================================================================

Let $\cA$ be a chiral algebra on a smooth projective curve~$X$,
equipped with its PBW filtration $F_\bullet$.
Write $\barBch(\cA) = T^c(s^{-1}\bar\cA)$ for the geometric chiral
bar complex with deconcatenation coproduct, $R_\cA = \gr^F \cA$ for
the associated graded (a Poisson vertex algebra), and
$\Omega^{\mathrm{ch}}$ for the cobar functor.

\begin{definition}[Chiral Koszul morphism]
\label{def:kfc-chiral-koszul-morphism}
A chiral twisting datum $(\cA, \cC, \tau, F_\bullet)$
consisting of an augmented chiral algebra~$\cA$, a conilpotent
factorization coalgebra~$\cC$, a degree-$+1$ Maurer--Cartan element
$\tau \colon \cC \to \cA$ in the convolution algebra, and a
compatible PBW filtration, is \emph{Koszul} if the following three
conditions hold:
\begin{enumerate}[label=\textup{(\alph*)}]
\item $\gr_F\cA$ is classically Koszul on each finite PBW window;
\item both twisted tensor products
\[
K_\tau^L(\cA,\cC) = (\cA \otimes \cC,\; d_\cA + d_\cC + d_\tau^L),
\qquad
K_\tau^R(\cC,\cA) = (\cC \otimes \cA,\; d_\cC + d_\cA + d_\tau^R)
\]
are acyclic;
\item the PBW spectral sequence is complete, separated, and
Mittag--Leffler on finite windows, so the associated graded
acyclicity lifts to the completed bar complex.
\end{enumerate}
The chiral algebra~$\cA$ is \emph{chirally Koszul}
if the canonical bar twisting datum
$(\cA, \barBch(\cA), \tau_0, F_\bullet^{\mathrm{PBW}})$ is Koszul.
\end{definition}

\begin{theorem}[Conditional chart criteria for chiral Koszulness]
\label{thm:kfc-fourteen}
Let $\cA$ be a positive-energy chiral algebra on~$X$ with finite
PBW windows, PBW-flat associated graded, and a complete separated
Mittag--Leffler PBW filtration. Conditions
\textup{(i)--(iv)}, \textup{(vi)}, \textup{(viii)--(x)}, and
\textup{(xii)--(xiii)} below are bidirectionally equivalent under
the Theorem-B ambient hypotheses for the chiral bar--cobar
adjunction.
Condition~\textup{(vii)} is a sufficient PBW-free chart:
flat PBW deformation with classically quadratic-Koszul associated
graded implies the core conditions. Freely strongly generated
vertex algebras satisfy this chart. The converse is not asserted
without an explicit PBW-free detection theorem.
Condition~\textup{(xi)} is equivalent to them for all chiral
algebras whose FM dual complex is simplicial.
Condition~\textup{(v)} is a proved one-way consequence on the
Koszul locus, with a conditional converse stated in
Theorem~\ref{thm:kfc-v-converse}.
Condition~\textup{(xiv)} is equivalent to the others for the
affine Kac--Moody family and one-directional in general.

Under the perfectness hypothesis on the ambient tangent complex,
the Lagrangian criterion~\textup{(xv)} is also equivalent.
$\cD$-module purity~\textup{(xvi)} implies~\textup{(x)}; the
converse is proved for affine Kac--Moody and open in general.

\medskip
\noindent
\textbf{Core equivalences}
\textup{(homological algebra and factorization geometry):}
\begin{enumerate}[label=\textup{(\roman*)},leftmargin=2.4em]

\item\label{item:kfc-koszul}
\textbf{Chiral Koszulness.}
The twisted tensor products $K_\tau^L(\cA, \barBch(\cA))$ and
$K_\tau^R(\barBch(\cA), \cA)$ are acyclic
\textup{(}Definition~\textup{\ref{def:kfc-chiral-koszul-morphism})}.

\item\label{item:kfc-pbw}
\textbf{PBW $\Etwo$-collapse.}
The PBW spectral sequence
$E_1^{p,q} = H^{p+q}(\gr^F_p \barBch(\cA))
 \Rightarrow H^{p+q}(\barBch(\cA))$
collapses at~$\Etwo$: all differentials $d_r = 0$ for $r \geq 2$.

\item\label{item:kfc-ainfty}
\textbf{$\Ainfty$-formality of bar cohomology.}
The minimal $\Ainfty$-model of $H^*(\barBch(\cA))$
is formal: the transferred higher products satisfy
$m_n = 0$ for all $n \geq 3$.

\item\label{item:kfc-shadow}
\textbf{Shadow-tower formality.}
The genus-zero obstruction classes of the shadow tower vanish:
$o_{r+1}^{(0)}(\cA) = 0$ for all $r \geq 2$.
Equivalently, the transferred $\Linfty$-brackets
$\ell_{r+1}^{(0)}$ of the modular convolution algebra vanish at
genus zero for all $r \geq 2$.

\item\label{item:kfc-e2formality}
\textbf{$\Etwo$-formality of chiral Hochschild cohomology.}
$\ChirHoch^*(\cA)$ is concentrated in cohomological degrees
$\{0,1,2\}$, carries a polynomial Hilbert series
$P_\cA(t) = \dim Z(\cA) + \dim \ChirHoch^1(\cA) \cdot t
 + \dim Z(\cA^!) \cdot t^2$,
and is formal as an $\Etwo$-algebra.

\item\label{item:kfc-curve-indep}
\textbf{Curve independence.}
The genus-zero factorization homology
$H^k(\int_{\bP^1}\cA)$ is independent of the choice of smooth
projective curve: replacing $\bP^1$ by any smooth proper
curve~$X'$ of genus zero gives the same cohomology.

\item\label{item:kfc-pbw-univ}
\textbf{PBW universality.}
The PBW filtration on~$\cA$ deforms flatly:
$\gr^F \cA_\hbar \cong \cA_0$ for the one-parameter deformation,
and the associated graded is quadratic-Koszul in the classical
sense.
This is a sufficient PBW-free condition, not an unconditional
equivalent of chiral Koszulness. Freely strongly generated vertex
algebras satisfy it. Simple quotients, logarithmic triplet
algebras, and symplectic-fermion orbifolds require separate
bar, determinant, or orbifold input; the quotient or orbifold
functor does not supply free strong generation by itself.

\item\label{item:kfc-bbl}
\textbf{Barr--Beck--Lurie monadicity.}
The comparison functor~$\Phi$ for the bar--cobar adjunction
$\barBch \dashv \Omega^{\mathrm{ch}}$ is an equivalence on the
fiber over $\barBch(\cA)$.

\item\label{item:kfc-fh}
\textbf{Factorization homology concentration.}
$H^k(\int_{\bP^1}\cA) = 0$ for $k \neq 0$.
Under the uniform-weight hypothesis, for every smooth projective
curve~$\Sigma_g$ of genus~$g \geq 0$,
$H^k(\int_{\Sigma_g}\cA) = 0$ for $k \neq 0$,
and the surviving scalar class at $g \geq 1$ is
$\mathrm{obs}_g(\cA) = \kappa(\cA) \cdot \lambda_g$
\textup{(UNIFORM-WEIGHT)}.

\item\label{item:kfc-fm}
\textbf{FM boundary acyclicity.}
For every Fulton--MacPherson boundary stratum
$S_T \cong \prod_{v \in V(T)} \FM_{|v|}(X)$
indexed by a rooted tree~$T$, the restriction
$i_{S_T}^!\,\barBch_n(\cA)$ is acyclic outside degree~$0$.

\end{enumerate}

\medskip
\noindent
\textbf{Tropical and representation-theoretic:}
\begin{enumerate}[label=\textup{(\roman*)},leftmargin=2.4em,start=11]

\item\label{item:kfc-tropical}
\textbf{Tropical Koszulness.}
The tropicalization of the bar complex
$\mathrm{trop}(\barBch(\cA))$, obtained by
passing to the dual complex of the FM compactification and
retaining only the combinatorial (tree-level) data, satisfies
tropical $\Etwo$-collapse:
the tropical spectral sequence has $d_r^{\mathrm{trop}} = 0$
for $r \geq 2$.

\item\label{item:kfc-ext}
\textbf{Ext diagonal vanishing.}
$\Ext^{p,q}_\cA(\omega_X, \omega_X) = 0$ for $p \neq q$,
where the bigrading comes from the PBW filtration on the bar
resolution.

\item\label{item:kfc-ks}
\textbf{Kac--Shapovalov determinant nonvanishing.}
The Shapovalov determinant $\det G_h \neq 0$ in the bar-relevant
range: the PBW-to-bar comparison map is injective at every
bar-relevant conformal weight~$h$.

\end{enumerate}

\medskip
\noindent
\textbf{Poisson-geometric (affine family):}
\begin{enumerate}[label=\textup{(\roman*)},leftmargin=2.4em,start=14]

\item\label{item:kfc-sklyanin}
\textbf{Sklyanin $H^2 = 0$.}
For $\cA = V_k(\fg)$ with $\fg$ semisimple, the second Poisson
cohomology of the Sklyanin bracket
$\{-,-\}_{\mathrm{STS}}$ on $(\fg^!)^*$ vanishes:
$H^2_\pi((\fg^!)^*, \{-,-\}_{\mathrm{STS}}) = 0$.
Equivalently, the Sklyanin--Poisson structure determined by the
classical $r$-matrix $r(z) = k\Omega/z$ admits no nontrivial
infinitesimal deformations.

\end{enumerate}

\medskip
\noindent
\textbf{Conditional and partial:}
\begin{enumerate}[label=\textup{(\roman*)},leftmargin=2.4em,start=15]

\item\label{item:kfc-lagrangian}
\textbf{Lagrangian criterion} \textup{(conditional on
perfectness).}
$\cM_\cA$ and $\cM_{\cA^!}$ are transverse Lagrangians in the
$(-1)$-shifted symplectic deformation space~$\cM_{\mathrm{comp}}$.

\item\label{item:kfc-purity}
\textbf{$\cD$-module purity} \textup{(one-directional).}
Each $\barBch_n(\cA)$ is pure of weight~$n$ as a mixed Hodge
module, with characteristic variety aligned to FM boundary strata.
Implies~\textup{\ref{item:kfc-fm}}. Converse proved for affine
Kac--Moody; open in general.

\end{enumerate}
\end{theorem}

\begin{remark}[Numbering convention]
\label{rem:kfc-numbering}
The numbering (i)--(xiv) groups conditions by
mathematical kinship.
The parent monograph~\cite{LorgatMKD1} numbers twelve
conditions (i)--(xii) in a slightly different order; the
present list expands by introducing shadow-tower
formality~(iv), $\Etwo$-formality~(v), curve
independence~(vi), PBW universality~(vii), and tropical
Koszulness~(xi) as named conditions, and by separating the
Sklyanin criterion~(xiv) from the Ext diagonal
vanishing~(xii).
The Lagrangian criterion~(xv) and $\cD$-module purity~(xvi)
are kept outside the numbered list because their equivalence
class is incomplete.
\end{remark}

\begin{remark}[Logical strength of the chart list]
\label{rem:fourteen-characterizations-strength}
The headline count \emph{fourteen} aggregates conditions of three
logical strengths; the list does not assert that all sixteen
labeled items~(i)--(xvi) are unconditional biconditionals.

\textbf{Unconditional biconditionals} (ten total):
(i)~Koszul, (ii)~PBW~$\Etwo$, (iii)~$\Ainfty$-formality,
(iv)~shadow-tower formality, (vi)~curve independence,
(viii)~Barr--Beck--Lurie monadicity,
(ix)~factorization-homology concentration,
(x)~FM boundary acyclicity, (xii)~$\Ext$ diagonal,
(xiii)~Kac--Shapovalov determinant.

\textbf{One-directional implications} (three total):
(v)~$\Etwo$-formality of $\ChirHoch^*$: direction
(i)$\Rightarrow$(v) proved, converse conditional in
Theorem~\ref{thm:kfc-v-converse}.
(vii)~PBW universality: the PBW-free condition implies the core
Koszul conditions; the converse would require a separate
detection theorem for free strong generation or for a
classically Koszul associated graded.
(xvi)~$\cD$-module purity: (xvi)$\Rightarrow$(x) proved,
converse for affine Kac--Moody only.

\textbf{Conditional on hypothesis} (two total):
(xi)~tropical Koszulness, equivalent when the FM dual complex
is simplicial (true for smooth curves).
(xv)~Lagrangian criterion, equivalent under perfectness and
non-degeneracy of the ambient tangent complex.

\textbf{Family-restricted} (one total):
(xiv)~Sklyanin $H^2 = 0$, biconditional for the affine
Kac--Moody family, one-directional ($\Rightarrow$(ii)) in
general.

Count: $10 + 3 + 2 + 1 = 16$ labeled items. The headline
\emph{fourteen} corresponds to (i)--(xiv); the conditional
pair (xv), (xvi) is kept outside.
\end{remark}

\begin{theorem}[Conditional Hochschild converse]
\label{thm:kfc-v-converse}
Let $\cA$ be a positive-energy chiral algebra on~$X$ with PBW
filtration. Suppose condition~(v) holds in its full form:
$\ChirHoch^*(\cA)$ is concentrated in cohomological degrees
$\{0,1,2\}$, carries the polynomial Hilbert series
$P_\cA(t) = \dim Z(\cA) + \dim \ChirHoch^1(\cA)\cdot t
 + \dim Z(\cA^!)\cdot t^2$,
and is formal as an $\Etwo$-algebra. Assume in addition:
\begin{enumerate}[label=\textup{(\alph*)}]
\item a two-colour chiral Higher Deligne comparison identifies
  $\ChirHoch^*(\cA)$ with the derived $\Etwo$-chiral centre of
  $\cA$ at the chosen associator;
\item the Koszul dual of \(\SCchtop\) admits a contracting
  homotopy on the finite PBW component under consideration;
\item the chiral cap product detects the bar cohomology
  operations faithfully.
\end{enumerate}
Then condition~\textup{(v)} implies chiral Koszulness
\textup{(i)}.
\end{theorem}

\begin{proof}
Four steps.

\emph{Step 1: Chiral Deligne-Tamarkin identification.} Under
the assumed two-colour chiral Higher Deligne comparison, the
$\Etwo$-algebra structure on $\ChirHoch^*(\cA)$ is
identified with the derived $\Etwo$-chiral centre
$\cZ^{\mathrm{der}}_{\mathrm{ch},\Etwo}(\cA)$ via
\[
\ChirHoch^*(\cA) \;\simeq\;
\RHom_{\Etwo^{\mathrm{ch}}\text{-}\mathrm{mod}}(\cA, \cA),
\]
where the right-hand side is computed in the factorization
ambient on~$\bP^1$. This identification is associator-dependent;
fix a Drinfeld associator~$\Phi$ (e.g.\ $\Phi_{\mathrm{KZ}}$).

\emph{Step 2: Cobar of formal $\Etwo$-algebra.} The
$\Etwo$-formality hypothesis on $\ChirHoch^*(\cA)$ gives a
strict $\Etwo$-algebra model. Kontsevich--Tamarkin operadic
Koszul duality between~$\Etwo$ and its shifted dual
$\Etwo\{1\}$ produces an $\Etwo$-cobar
$\Omega_{\Etwo}(\ChirHoch^*(\cA))$ that computes
$\RHom_{\Etwo\text{-}\mathrm{mod}}(k, \cA)$ in the chiral
ambient. The polynomial Hilbert series hypothesis
($\dim = \dim Z(\cA) + \dim\ChirHoch^1\cdot t
 + \dim Z(\cA^!)\cdot t^2$, no higher) forces the brace
operations $\mathrm{br}_n$ on $\ChirHoch^*(\cA)$ to vanish for
all $n \geq 3$: any nonzero $\mathrm{br}_n$ would contribute a
class in degree $\geq 3$ by arity-counting on
$\FM_n(\bC)$-chains, contradicting concentration in $\{0,1,2\}$.

\emph{Step 3: Transfer to $\Ainfty$-operations on bar
cohomology.} The bar coalgebra $\barBch(\cA)$ and its dual
$\RHom(\barBch(\cA), k) = C^*_{\mathrm{ch}}(\cA, k)$ are
related to $\ChirHoch^*(\cA)$ by the chiral cap product. Under
the chiral Deligne identification of Step~1, the transferred
$\Ainfty$-operations $m_n$ on $H^*(\barBch(\cA))$ are images of
the brace operations $\mathrm{br}_n$ under the Koszul dual
pairing. Brace vanishing for $n \geq 3$ (Step~2) forces
$m_n = 0$ for all $n \geq 3$. This is condition~(iii),
$\Ainfty$-formality of bar cohomology.

\emph{Step 4: Apply the core criterion.} By the core circuit of
Theorem~\ref{thm:kfc-fourteen}, (iii)~$\Leftrightarrow$~(i):
$\Ainfty$-formality of bar cohomology is equivalent to chiral
Koszulness. Hence (v)$\Rightarrow$(iii)$\Rightarrow$(i) under the
three comparison hypotheses.

The proof requires the associator choice~$\Phi$ of Step~1, so
the converse is chart-dependent. Associator-independent
two-colour formality is a separate theorem-level obligation.
\end{proof}

\begin{remark}[Dependence on chiral Higher Deligne]
\label{rem:kfc-v-converse-dependence}
The theorem isolates the exact hypotheses needed for the converse:
two-colour chiral Higher Deligne, a finite-PBW contracting homotopy
for the Koszul dual Swiss-cheese complex, and faithful cap-product
detection. The associator-independent formulation remains open.
\end{remark}

\begin{remark}[The four objects]
\label{rem:kfc-four-objects}
The four objects of chiral Koszul duality must not be conflated:
\begin{enumerate}
\item $\barBch(\cA) = T^c(s^{-1}\bar\cA)$: the bar
  \emph{coalgebra}, with deconcatenation coproduct and bar
  differential.
\item $\cA^i = H^*(\barBch(\cA))$: the dual
  \emph{coalgebra}, Koszul cohomology of the bar complex.
\item $\cA^! = ((\cA^i)^\vee)$: the Koszul dual
  \emph{algebra}, obtained by Verdier duality.
\item $\cZ^{\mathrm{der}}_{\mathrm{ch}}(\cA) =
  \RHom_{\cA \otimes \cA^{\mathrm{op}}}(\cA,\cA)$: the
  derived chiral center, computed using the bar complex as a
  resolution.
\end{enumerate}
The bar complex is the $\Eone$ engine; $\Omega(\barBch(\cA)) = \cA$
is \emph{inversion}, not Koszul duality. The Koszul dual comes from
Verdier duality. The derived center comes from chiral Hochschild cochains.
\end{remark}


% ================================================================
\section{The circuit of implications}
\label{sec:circuit}
% ================================================================

The proof of Theorem~\ref{thm:kfc-fourteen} is organized as a
directed graph of implications. The central structure is the
\emph{core circuit}, a cycle of biconditional implications among five
conditions:
\begin{equation}\label{eq:core-circuit}
\mathrm{(i)} \;\Leftrightarrow\;
\mathrm{(ii)} \;\Leftrightarrow\;
\mathrm{(iii)} \;\Leftrightarrow\;
\mathrm{(viii)} \;\Leftrightarrow\;
\mathrm{(ix)}.
\end{equation}
Each of the remaining conditions is related to the core by the
satellite implication shown in the graph; bidirectional arrows
are equivalences, while single arrows are one-way criteria.
The full graph has the following shape.

\begin{center}
\begin{tikzcd}[row sep=1.6em, column sep=1.8em]
& \text{(iv) shadow form.} \ar[d,Leftrightarrow] \\
\text{(xii) Ext diag.} \ar[r,Leftrightarrow]
& \text{(ii) PBW $E_2$} \ar[r,Leftrightarrow]
 \ar[d,Leftrightarrow]
& \text{(i) Koszul} \ar[r,Leftrightarrow]
& \text{(viii) BBL} \\
\text{(xiii) Kac--Shap.} \ar[ur,Leftrightarrow]
& \text{(iii) $A_\infty$ form.} \\
\text{(vi) curve indep.} \ar[uur,Leftrightarrow]
& \text{(vii) PBW univ.}
 \ar[u,Rightarrow] \ar[uu,Rightarrow,crossing over] \\
\text{(x) FM bdry} \ar[uuur,Leftrightarrow]
& \text{(xi) tropical} \ar[l,Leftrightarrow] \\
\text{(v) $E_2$-Hoch} & \text{(ix) FH conc.}
 \ar[uuuu,Leftrightarrow,bend right=40] \\
\text{(xiv) Skl.\ $H^2$} \ar[uuuu,Rightarrow]
& \text{(xv) Lagr.} \ar[uuuuu,Leftrightarrow,dashed]
& \text{(xvi) purity} \ar[uuuul,Rightarrow,dashed]
\end{tikzcd}
\end{center}

\noindent
Single arrows indicate one-way implications; dashed arrows indicate
conditional or partial implications.
We now give proof sketches for each link, following the order
of the graph from the core outward.

\subsection{The core circuit}

\subsubsection*{\textup{(i)}~$\Leftrightarrow$~\textup{(ii):}
Koszulness and PBW collapse}

Koszulness is acyclicity of $K_\tau^L(\cA, \barBch(\cA))$
for the canonical twisting morphism~$\tau_0$. The left
twisted tensor product is filtered by the PBW bar degree, and on
the associated graded the twisted complex becomes the standard
Koszul complex of the Poisson vertex algebra~$R_\cA$.
Acyclicity on the associated graded lifts to acyclicity of
$K_\tau^L$ if and only if the spectral sequence collapses
at~$\Etwo$~\cite[Thm.~4.2.1]{LorgatMKD1}. This is the PBW
Koszulness criterion.

Conversely, $\Etwo$-collapse means $H^*(\barBch(\cA))$ is
concentrated on the PBW diagonal $p = q$. A $d_r$-differential
in the spectral sequence moves off-diagonal by $(r, 1-r)$;
off-diagonal classes in the abutment would contradict diagonal
concentration. The spectral sequence converges conditionally in
the sense of Boardman (the PBW filtration on a positive-energy
chiral algebra is bounded below and exhaustive), and in
characteristic zero with a split filtration the extension problem
is trivial.

\subsubsection*{\textup{(ii)}~$\Leftrightarrow$~\textup{(iii):}
PBW collapse and $\Ainfty$-formality}

The Homological Perturbation Lemma identifies the differentials
$d_r$ of the PBW spectral sequence with the transferred
$\Ainfty$-operations $m_{r+1}$ on bar cohomology, up to sign.
Hence $d_r = 0$ for all $r \geq 2$ is equivalent to $m_n = 0$
for all $n \geq 3$: the minimal $\Ainfty$-model is formal.

For the converse, if $m_n = 0$ for $n \geq 3$, the minimal model
is a strict dg algebra. The chiral $\Ainfty$-structure is computed
fiberwise on each FM stratum; on each fiber, Keller's classicality
theorem for formal $\Ainfty$-algebras gives $\Etwo$-collapse.
The PBW filtration is defined fiberwise and compatible with the FM
stratification, so fiberwise collapse assembles to global collapse.

\subsubsection*{\textup{(i)}~$\Leftrightarrow$~\textup{(viii):}
Koszulness and Barr--Beck--Lurie monadicity}

The Barr--Beck--Lurie theorem~\cite[Thm.~4.7.3.5]{HA} states that
the comparison functor~$\Phi$ for the bar--cobar adjunction is an
equivalence if and only if the bar functor is conservative and
preserves totalizations of bar-split cosimplicial objects.

On the Koszul locus, the bar functor detects quasi-isomorphisms
(the bar--cobar counit inverts them), giving conservativity. The
bar--cobar adjunction is a Quillen equivalence on the Koszul
locus~\cite[\S5]{LorgatMKD1}; a Quillen equivalence of stable
model categories satisfies the monadic descent conditions of
\cite[Prop.~4.7.3.16]{HA}. Hence BBL monadicity holds if and
only if the chiral algebra is Koszul.

\subsubsection*{\textup{(i)}~$\Leftrightarrow$~\textup{(ix):}
Koszulness and factorization homology concentration}

The bar complex computes factorization homology:
$\barBch_n(\cA) = (\int_X \cA)_n$. At genus zero,
$\Etwo$-collapse gives
$H^k(\int_{\bP^1}\cA) = 0$ for $k \neq 0$: the Goodwillie
filtration has associated graded given by bar contributions,
each concentrated in degree~$0$ by PBW collapse.
Under the uniform-weight hypothesis, the same holds at every
genus: the loop-order spectral sequence collapses on the
modular Koszul locus, each vertex contribution remains in
degree~$0$, and the surviving scalar is
$\mathrm{obs}_g(\cA) = \kappa(\cA) \cdot \lambda_g$
(UNIFORM-WEIGHT).

Conversely, concentration at genus zero specializes to
$H^k(\int_{\bP^1}\cA) = 0$ for $k \neq 0$. This is computed
by the bar complex: the geometric realization
$\int_{\bP^1}\cA \simeq |\barBch(\cA)|$. Concentration in
degree~$0$ forces the bar spectral sequence to collapse, which
by the first equivalence is Koszulness.

\subsection{Satellite equivalences from the core}

\subsubsection*{\textup{(ii)}~$\Leftrightarrow$~\textup{(iv):}
PBW collapse and shadow-tower formality}

The shadow obstruction tower $\{o_{r+1}(\cA)\}_{r \geq 2}$
has two projections: the genus-zero slice
$o_{r+1}^{(0)}(\cA)$ and the all-genera projection. Shadow-tower
formality (vanishing of the genus-zero slice) is identified with the
transferred $\Linfty$-brackets $\ell_{r+1}^{(0)}$ of the modular
convolution algebra at genus zero. The HPL transfer identifies
these with the PBW spectral sequence differentials.
Hence shadow-tower formality is equivalent to $\Etwo$-collapse.

The all-genera shadow tower, by contrast, can be nontrivial even for
Koszul algebras: this is the shadow depth classification G/L/C/M.
Shadow-tower formality is the genus-zero condition; shadow depth
is the all-genera condition. They are independent invariants.

\subsubsection*{\textup{(i)}~$\Rightarrow$~\textup{(v):}
Koszulness implies $\Etwo$-formality of ChirHoch}

The bar--cobar quasi-isomorphism
$\Omega^{\mathrm{ch}}(\barBch(\cA)) \xrightarrow{\sim} \cA$
provides a free resolution of~$\cA$. Computing
$\ChirHoch^*(\cA)$ through this resolution, the bar--cobar
filtration forces diagonal concentration. The polynomial
structure in degrees $\{0,1,2\}$ follows: degree~$0$ is the
center, degree~$1$ is first-order deformations, degree~$2$ is
obstructions. The $\Etwo$-formality of $\ChirHoch^*(\cA)$ uses
the formality of the local configuration spaces $\FM_n(\bC)$.

This is a one-way consequence: $\Etwo$-formality of $\ChirHoch^*$
does not force the underlying graded-commutative algebra to be free,
so the converse from (v) to bar--cobar inversion is not established.

\subsubsection*{\textup{(ii)}~$\Leftrightarrow$~\textup{(vi):}
PBW collapse and curve independence}

Curve independence at genus zero follows from PBW collapse because
the collapsed spectral sequence depends only on the associated
graded $R_\cA$, which is a Poisson vertex algebra independent of
the global geometry of~$X$. The FM compactification
$\FM_n(X)$ depends on~$X$, but the $\Etwo$-page
$H^{p+q}(\gr^F_p \barBch(\cA))$ is computed from $R_\cA$ by
Chevalley--Eilenberg methods that see only the local OPE data.

Conversely, if factorization homology is curve-independent at
genus zero, the bar complex cannot have differentials sensitive to
global geometry. The PBW spectral sequence differentials
$d_r$ for $r \geq 2$ involve integration over FM strata of
codimension~$r$; these strata encode collisions of $r+1$ points
and are sensitive to the curve's geometry at the $r$-fold
collision scale. Curve independence forces these to vanish,
giving $\Etwo$-collapse.

\subsubsection*{\textup{(vii)}~$\Rightarrow$~\textup{(iii):}
PBW-free universality implies $\Ainfty$-formality}

PBW universality states that the PBW filtration deforms flatly
with classically Koszul associated graded. If $\cA$ is freely
strongly generated (the generators form a free $\cD_X$-module),
the associated graded is a free polynomial PVA, which is
quadratic-Koszul. The flat deformation property then implies that
the bar spectral sequence for $\cA_\hbar$ is a deformation of
the (collapsed) spectral sequence for $\cA_0$. By
semicontinuity, the collapsed page persists, giving
$\Etwo$-collapse and hence $\Ainfty$-formality.

This proves the satellite implication
\textup{(vii)}$\Rightarrow$\textup{(iii)}. The reverse implication
is not used here: $\Ainfty$-formality of bar cohomology does not
by itself reconstruct a free strong generating set, a polynomial
PVA presentation, or a classically Koszul associated graded.

\subsubsection*{\textup{(i)}~$\Leftrightarrow$~\textup{(x):}
Koszulness and FM boundary acyclicity}

Each FM boundary stratum $S_T$ is indexed by a rooted tree~$T$.
The residue component $d_{\mathrm{res}}|_{S_T}$ of the bar
differential restricts to the tensor product of lower-degree bar
complexes at the vertices of~$T$. $\Etwo$-collapse at each
vertex gives acyclicity at $S_T$:
$H^k(i_{S_T}^!\,\barBch_n(\cA)) = 0$ for $k \neq 0$.

Conversely, acyclicity at the deepest strata (binary collision
divisors $S = D_{\{i,j\}} \cong X \times \FM_{n-1}(X)$) forces
every OPE residue to be exact in the PBW-graded sense, which is
the condition for $\Etwo$-collapse. Iterating over all strata
closes the equivalence.

\subsubsection*{\textup{(x)}~$\Leftrightarrow$~\textup{(xi):}
FM boundary acyclicity and tropical Koszulness}

The tropicalization of the bar complex retains only the
combinatorial data visible in the dual complex of the FM
compactification: the rooted trees indexing boundary strata, the
face maps between strata, and the residues of the bar differential
along each stratum. Tropical $\Etwo$-collapse is the assertion
that these combinatorial residues produce no off-diagonal cohomology.

FM boundary acyclicity (the analytic statement) implies tropical
Koszulness (the combinatorial statement) by restriction. The
converse holds when the dual complex is simplicial (the standard
case for smooth curves): the combinatorial data determines the
analytic behavior up to homotopy equivalence, so tropical acyclicity
forces analytic acyclicity.

\subsubsection*{\textup{(i)}~$\Leftrightarrow$~\textup{(xii):}
Koszulness and Ext diagonal vanishing}

$\Etwo$-collapse places $H^*(\barBch(\cA))$ on the PBW diagonal
$p = q$. The bar complex computes
$\Ext^{p,q}_\cA(\omega_X, \omega_X)$, so diagonal concentration
of bar cohomology is Ext diagonal vanishing.

Conversely, Ext diagonal vanishing means the abutment of the PBW
spectral sequence is diagonal. Any nonzero differential $d_r$ for
$r \geq 2$ would produce off-diagonal classes, contradicting
diagonal concentration. Hence $d_r = 0$ for all $r \geq 2$.

\subsubsection*{\textup{(i)}~$\Leftrightarrow$~\textup{(xiii):}
Koszulness and Kac--Shapovalov determinant}

If $\cA$ is Koszul, the PBW filtration is strict on the vacuum
module: $\gr^F M_0(\cA) \cong \Symch(V)$. Strictness implies
non-degeneracy of the Shapovalov form $G_h$ in the bar-relevant
range.

Conversely, non-degeneracy of $G_h$ makes the PBW-to-bar comparison
injective at every bar-relevant weight. Suppose $d_r \neq 0$ for
some $r \geq 2$; then there exists a nonzero class in
$E_r^{p,q}$ with nonzero differential, whose lift to the bar
complex contradicts the acyclicity of $\Sym(V)$ as a Koszul
complex. Hence $d_r = 0$ for all $r \geq 2$.

\subsubsection*{\textup{(xiv)}~$\Rightarrow$~\textup{(ii):}
Sklyanin $H^2 = 0$ implies PBW collapse (affine family)}

For $\cA = V_k(\fg)$ with $\fg$ semisimple, the Koszul dual
coalgebra $\barBch(\cA)^!$ carries the Poisson-Lie group
$(\fg^!)^*$ with the Sklyanin--Poisson bracket determined by the
classical $r$-matrix $r(z) = k\Omega/z$ via the
Semenov-Tian-Shansky construction. The second Poisson cohomology
$H^2_\pi$ controls infinitesimal deformations of this bracket.

The vanishing $H^2_\pi = 0$ is established by Whitehead's second
lemma (for $\fg$ semisimple, $H^2(\fg; M) = 0$ for finite-dimensional
modules~$M$): the Poisson cochain complex
$(\mathfrak{X}^\bullet, \delta_\pi)$ is the associated graded of
the bar differential under PBW filtration, so vanishing of $H^2_\pi$
is the associated-graded avatar of $\Etwo$-collapse.
The proof: five independent derivation paths (direct kernel/image,
Chevalley--Eilenberg comparison, polynomial-degree spectral sequence,
Whitehead reduction, and numerical from the compute engine) confirm
$H^0_\pi = \bC[C_2]$ (Casimir ring),
$H^1_\pi = H^2_\pi = H^3_\pi = 0$.

This is proved as a biconditional for the affine Kac--Moody family.
For general chiral algebras, the Sklyanin construction does not
apply (it requires a Lie-algebraic $r$-matrix), so the equivalence is
restricted to this family.

\subsubsection*{\textup{(xv):}
Lagrangian criterion, conditionally}

Under the perfectness and non-degeneracy hypotheses on the ambient
tangent complex: on the Koszul locus, the complementarity splitting
\[
\mathbf{Q}_g(\cA) \oplus \mathbf{Q}_g(\cA^!)
 \;\simeq\; H^*(\overline{\cM}_g, \cZ_\cA)
\]
identifies $\cM_\cA$ and $\cM_{\cA^!}$ as complementary subspaces
of $\cM_{\mathrm{comp}}$. The $(-1)$-shifted symplectic structure
makes each isotropic; complementarity gives maximal dimension,
hence Lagrangian. Conversely, transverse Lagrangian intersection
in a $(-1)$-shifted symplectic space has expected dimension zero;
its discreteness is acyclicity of the twisted tensor product,
which is Koszulness.

\subsubsection*{\textup{(xvi)}~$\Rightarrow$~\textup{(x):}
$\cD$-module purity implies FM boundary acyclicity}

Suppose each $\barBch_n(\cA)$ is pure of weight~$n$ as a mixed
Hodge module. By Saito's theory, the weight filtration on a pure
Hodge module splits up to semisimplification; purity plus the
alignment of the characteristic variety to FM strata forces
acyclicity outside degree~$0$ on every stratum. The converse
(Koszulness implies purity) is proved for the affine Kac--Moody
family, where the bar complex is identified with a geometric
complex whose purity follows from the Riemann hypothesis for
weight filtrations.

% ================================================================
\section{Scalar saturation}
\label{sec:scalar-saturation}
% ================================================================

\subsection{The modular characteristic}

Let $\cA$ be a chiral algebra on~$X$. The universal Maurer--Cartan
element $\Theta_\cA := D_\cA - d_0$ of the bar differential lives
in the modular convolution dg Lie algebra
$\fg^{\mathrm{mod}}_\cA$ and satisfies
\begin{equation}\label{eq:mc-equation}
d\,\Theta_\cA + \tfrac{1}{2}[\Theta_\cA, \Theta_\cA] = 0.
\end{equation}
At each degree $r \geq 2$, $\Theta_\cA$ projects to a shadow
invariant $S_r(\cA)$.

The $\Eone$ ordered bar complex
$\barB^{\mathrm{ord}}(\cA) = T^c(s^{-1}\bar\cA)$ is the primitive
object: it carries the full $R$-matrix
$r_\cA(z) = \mathrm{Res}^{\mathrm{coll}}_{0,2}(\Theta_\cA)$.
The symmetric bar $\barB^\Sigma(\cA)$ is the
$\Sigma_n$-coinvariant quotient: a lossy projection that discards
the matrix-valued data and retains only the scalar
\begin{equation}\label{eq:kappa-definition}
\kappa(\cA) \;=\; S_2(\cA),
\end{equation}
the \emph{modular characteristic}: the single number that survives
averaging.

\begin{remark}[Family-specific formulas]
\label{rem:kfc-kappa-families}
The modular characteristic is family-specific; writing it from
the canonical landscape table prevents convention drift. The
verified formulas are:
\begin{center}
\renewcommand{\arraystretch}{1.3}
\begin{tabular}{lll}
\toprule
Family & $\kappa(\cA)$ & Sanity checks \\
\midrule
Heisenberg $\mathcal{H}_k$
  & $k$
  & $k=0 \to 0$; $k=1 \to 1$ \\
Affine KM $V_k(\fg)$
  & $\dfrac{\dim(\fg)(k+h^\vee)}{2h^\vee}$
  & $k=0 \to \frac{\dim(\fg)}{2}$;
    $k=-h^\vee \to 0$ \\
Virasoro $\Vir_c$
  & $c/2$
  & $c=0 \to 0$; $c=13 \to 13/2$ \\
$\cW_N$ (principal)
  & $c \cdot (H_N - 1)$
  & $N=2 \to c/2 = \kappa^{\Vir}$;
    $N=3 \to 5c/6$ \\
\bottomrule
\end{tabular}
\end{center}
\noindent
where $H_N = \sum_{j=1}^{N} 1/j$ is the $N$-th harmonic number.
\end{remark}

\subsection{The saturation theorem}

\begin{theorem}[Scalar saturation on finite PBW charts]
\label{thm:kfc-scalar-saturation}
Let $\cA_k$ be a one-parameter algebraic family of chiral algebras
with rational OPE coefficients, parametrized by a level~$k$.
Assume a finite PBW-Artin chart, scalar-diagonal Zamolodchikov
metric, and a one-dimensional degree-two deformation space.
Then the scalar projection of the degree-two universal MC element
is determined by the single scalar $\kappa(\cA_k)$:
\begin{equation}\label{eq:scalar-saturation}
\av\!\left(\Theta_{\cA_k}^{(2)}\right)
 \;=\; \kappa(\cA_k) \cdot \eta \otimes \Lambda,
\end{equation}
where $\eta = d\log(z_1 - z_2)$ is the Arnold propagator and
$\Lambda$ is the universal quadratic Casimir element.
\end{theorem}

\begin{proof}[Proof sketch]
The degree-two component of $\Theta_\cA$ in the modular convolution
algebra takes values in
$H^2(\fg_\cA^{\mathrm{mod}}; \Sym^q V)$ for appropriate~$q$.
Whitehead's second lemma gives vanishing of this cohomology for
reductive Lie algebras at generic level: the space of degree-two
MC deformations is one-dimensional, spanned by
$\eta \otimes \Lambda$. The coefficient is determined by the
$\Sigma_n$-coinvariant projection:
$\av(\Theta_\cA^{(2)}) = \kappa(\cA) \cdot [\lambda_1]$, which
fixes the coefficient as~$\kappa$.

The argument has three steps:
\begin{enumerate}
\item \textbf{Whitehead reduction.} For each standard family, the
  Lie algebra cohomology $H^2(\fg; \Sym^q V) = 0$ at generic level
  (the reductive hypothesis is used here). This gives a
  one-dimensional deformation space at degree two.
\item \textbf{Algebraic semicontinuity.} The exceptional locus
  where Whitehead reduction fails is Zariski-closed in the level
  parameter~$k$. The theorem is asserted on the complement of this
  exceptional locus and on the named standard rows where the
  finite-PBW computation has been carried out.
\item \textbf{Coefficient identification.} The averaging map
  $\av \colon \fg_\cA^{\Eone} \to \fg_\cA^{\mathrm{mod}}$
  projects $\Theta_\cA^{(2)}$ to
  $\kappa(\cA) \cdot [\lambda_1]$. Since $\av$ is the
  $\Sigma_n$-coinvariant projection and $\eta \otimes \Lambda$ is
  the unique $\Sigma_2$-invariant degree-two element (up to scalar),
  the coefficient is~$\kappa$.
\end{enumerate}

For non-abelian affine KM, the averaging identity involves the
Sugawara shift:
$\av(r(z)) + \dim(\fg)/2 = \kappa(V_k(\fg))$
where $\av(r(z)) = k \cdot \dim(\fg)/(2h^\vee)$ is the
double-pole channel and $\dim(\fg)/2$ is the simple-pole
self-contraction through the adjoint Casimir.
\end{proof}

\begin{remark}[Scope]
\label{rem:kfc-saturation-scope}
Scalar saturation at degree~$2$ is proved on the scalar-diagonal
standard rows and on finite PBW-Artin components satisfying the
hypotheses of Theorem~\ref{thm:kfc-scalar-saturation}. For
multi-weight algebras at genus $g \geq 2$, the scalar projection
saturates while the full MC element receives cross-channel
corrections. The stronger identity
$\Theta_\cA^{\mathrm{min}} = \kappa \cdot \eta \otimes \Lambda$
is proved under the uniform-weight hypothesis. At $g \geq 2$ it
receives cross-channel corrections for multi-weight families.
\end{remark}

\subsection{From scalar saturation to Koszulness}

Scalar saturation establishes that $\kappa$ controls the MC element
at degree~$2$. The question is what controls it at \emph{all}
degrees. The answer is the Koszul condition.

\begin{proposition}[Koszulness controls the full tower]
\label{prop:kfc-koszul-controls-tower}
If $\cA$ is chirally Koszul, then the universal MC element
$\Theta_\cA$ is determined at every degree by its genus-zero
projection. Concretely:
\begin{enumerate}
\item The shadow invariants $S_r(\cA)$ for $r \geq 3$ are
  determined by the $\Ainfty$-formality data of bar cohomology.
\item Under the uniform-weight hypothesis, $\Theta_\cA$ is determined
  by $\kappa(\cA)$ and the shadow depth class.
\item For class~G algebras (Heisenberg), $S_r = 0$ for
  $r \geq 3$: the MC element is entirely controlled by~$\kappa$.
\item For class~M algebras (Virasoro), $S_r \neq 0$ for
  infinitely many~$r$, but each $S_r$ is determined by
  $\kappa = c/2$ and $S_4 = 10/[c(5c+22)]$.
\end{enumerate}
\end{proposition}

\begin{proof}[Proof sketch]
$\Ainfty$-formality (condition~(iii) of
Theorem~\ref{thm:kfc-fourteen}) means $m_n = 0$ for $n \geq 3$:
the bar cohomology algebra has no higher operations. The MC moduli
of a formal $\Ainfty$-algebra is discrete; the shadow invariants
$S_r$ are the Taylor coefficients of the unique MC solution
in the completion of the convolution algebra. Each $S_r$ is
determined recursively from $\kappa = S_2$ and the structure
constants of the algebra.

For class~G, all $S_r = 0$ for $r \geq 3$ by direct computation
(the bar complex is abelian).
For class~M, the recursion is controlled by the Riccati equation
$\Delta = 8\kappa S_4$, where
$\Delta = 0$ characterizes class~G. The seed $S_4$ determines
$S_r$ for all $r$ via the shadow algebra Lie bracket; for
Virasoro, $S_4 = 10/[c(5c+22)] \sim 2/c^2$ at large~$c$.
\end{proof}

\begin{remark}[The critical discriminant]
\label{rem:kfc-discriminant}
The discriminant $\Delta = 8\kappa \cdot S_4$ is linear in
$\kappa$ (not quadratic). $\Delta = 0$ if and only if the shadow
tower is finite. For $\kappa \neq 0$, finiteness is equivalent to
$S_4 = 0$, which characterizes class~G. For $\kappa = 0$ (critical
level), $\Delta = 0$ regardless of~$S_4$.
\end{remark}

\subsection{The $\Eone$ structure}

Scalar saturation is a statement about the $\Sigma_n$-coinvariant
quotient. The full $\Eone$ ordered bar complex carries strictly
more data: the matrix-valued classical $r$-matrix
$r_\cA(z) \in \mathrm{End}_\cA(2)[\![z^{-1}]\!]$, of which
$\kappa$ is only the scalar shadow $\kappa = \av(r_\cA(z))$.

The $r$-matrix satisfies the classical Yang--Baxter equation
$[r_{12}, r_{13}] + [r_{12}, r_{23}] + [r_{13}, r_{23}] = 0$
and determines the Yangian structure on the ordered bar cohomology.
It absorbs one pole from the OPE via the logarithmic propagator:
for Heisenberg, the OPE has a double pole and the $r$-matrix
has a simple pole, $r^{\Heis}(z) = k/z$. For Virasoro, the OPE has
a quartic pole and the $r$-matrix has a cubic pole,
$r^{\Vir}(z) = (c/2)/z^3 + 2T/z$.
For affine KM in the trace-form convention,
$r^{\KM}(z) = k\Omega/z$, where $\Omega$ is the inverse Killing
form Casimir and the level prefix~$k$ is mandatory: at $k = 0$
the $r$-matrix vanishes (abelian limit).

The averaging map $\av \colon \fg_\cA^{\Eone} \to
\fg_\cA^{\mathrm{mod}}$ is the $\Sigma_n$-coinvariant projection.
It is lossy: the full $r$-matrix determines $\kappa$, but $\kappa$
does not determine the $r$-matrix. The five theorems of the parent
monograph extract the $\Sigma_n$-invariant content; the present
paper works at this invariant level, where scalar saturation gives
full control at degree two and Koszulness extends control to all
degrees.


% ================================================================
\section{The standard landscape}
\label{sec:landscape}
% ================================================================

The fourteen characterizations of
Theorem~\ref{thm:kfc-fourteen} are verified on the standard
generic rows of the landscape. The verification proceeds family
by family, each through a criterion adapted to the structure of
the algebra.

\begin{theorem}[Standard generic families]
\label{thm:kfc-landscape}
The following chiral algebras are chirally Koszul:
Heisenberg~$\mathcal{H}_k$ at all levels~$k$; lattice
algebras $V_L$ for all even lattices~$L$; affine Kac--Moody
$V_k(\fg)$ for semisimple~$\fg$ at generic non-critical
levels; Virasoro~$\Vir_c$ at generic central charge; the
principal $\cW$-algebra $\cW_N$ at all $N$ and generic level; the
$\beta\gamma$ system at standard conformal weight; and the
Bershadsky--Polyakov algebra~$\cW^{(2)}_k(\mathfrak{sl}_3)$ at
generic level.
\end{theorem}

The proof is the content of this section. Table~\ref{tab:landscape}
records the family, the Koszul equivalence that proves it, the
shadow class, the modular characteristic, and the Koszul
complementarity conductor.

\begin{table}[ht]
\centering
\renewcommand{\arraystretch}{1.3}
\begin{tabular}{lllll}
\toprule
Family & Koszul equivalence & Class & $\kappa$ & $K$ \\
\midrule
$\mathcal{H}_k$ & (ii) PBW, trivial & G & $k$ & $0$ \\
Lattice $V_L$ & (ii) PBW, lattice & G & $\mathrm{rk}(L)/2$ & $0$ \\
$V_k(\fg)$ & (xiii) Kac--Shapovalov & L & $\frac{\dim(\fg)(k+h^\vee)}{2h^\vee}$ & $0$ \\
$\Vir_c$ & (iv) shadow formality & M & $c/2$ & $13$ \\
$\cW_N$ & (xiii) via DS from KM & M & $c(H_N-1)$ & family-dep. \\
$\beta\gamma$ & (iii) $\Ainfty$-formality & C & $6\lambda^2-6\lambda+1$ & $0$ \\
$\cW^{(2)}_k(\mathfrak{sl}_3)$ & (xii) Ext diag. & M & $c(k)/6$ & $98/3$ \\
\bottomrule
\end{tabular}
\caption{The standard landscape and Koszulness.
  $H_N = \sum_{j=1}^{N} 1/j$. The column $K$ denotes the Koszul
  complementarity conductor $K = \kappa(\cA) + \kappa(\cA^!)$.}
\label{tab:landscape}
\end{table}

\subsection{Heisenberg}

The Heisenberg algebra $\mathcal{H}_k$ has a single generator
$J$ with OPE $J(z) J(w) \sim k/(z-w)^2$, so the bar complex
$\barBch(\mathcal{H}_k)$ is the desuspended tensor coalgebra on
a one-dimensional augmentation ideal. The PBW spectral sequence
has a single nonzero column at $p = 1$, so $\Etwo$-collapse
is vacuous: there are no possible higher differentials. This is
equivalence~(ii) of Theorem~\ref{thm:kfc-fourteen}, in its
trivial form. The modular characteristic is
$\kappa(\mathcal{H}_k) = k$, and the shadow tower terminates
at $r = 2$: class~G. The Koszul dual is
$\mathcal{H}_k^! = \Symch(V^*)$ (the chiral symmetric algebra
on the linear dual), with
$\kappa(\mathcal{H}_k^!) = -k$, giving
$K = k + (-k) = 0$.

\subsection{Lattice algebras}

The lattice vertex algebra $V_L$ associated to an even
lattice~$L$ of rank~$r$ is an extension of the Heisenberg
algebra by vertex operators $e^\alpha$ for $\alpha \in L$.
The PBW filtration on $V_L$ has associated graded
$\gr^F V_L \cong \Symch(\fh) \otimes \bC[L]$, where
$\fh = L \otimes_\bZ \bC$. The lattice grading gives a
decomposition of the bar complex into weight spaces, each of
which reduces to a Heisenberg bar complex tensored with
the group algebra. The PBW spectral sequence collapses at
$\Etwo$ in each weight sector: equivalence~(ii).
The modular characteristic is
$\kappa(V_L) = \mathrm{rk}(L)/2$, and the shadow class is~G.

\subsection{Affine Kac--Moody}

For $\cA = V_k(\fg)$ with $\fg$ semisimple and $k \neq -h^\vee$,
the Koszulness proof goes through the Kac--Shapovalov
criterion~(xiii). The Shapovalov form $G_h$ on the vacuum Verma
module $M_0(V_k(\fg))$ is nondegenerate for $k$ outside the
critical hyperplane: the Shapovalov determinant is a product
over positive roots and PBW-relevant weights, and its zeros are
confined to the critical level and the reducibility
locus~\cite[Thm.~11.4]{KK79}.

At generic non-critical level, $\det G_h \neq 0$ in the
bar-relevant range: the PBW-to-bar comparison is injective, and the
PBW spectral sequence collapses at~$\Etwo$ by
Theorem~\ref{thm:kfc-fourteen}~(xiii)$\Leftrightarrow$(ii).
This is the chiral analogue of the Garland--Lepowsky
concentration theorem~\cite{GL76}: the cohomology of the bar
complex is concentrated on the PBW diagonal, the higher
differentials vanish, and the associated graded computes the
Koszul dual coalgebra.

The modular characteristic is
$\kappa(V_k(\fg)) = \dim(\fg)(k+h^\vee)/(2h^\vee)$.
Sanity checks: $k = 0$ gives $\kappa = \dim(\fg)/2$ in the
abelian limit: the double-pole vanishes and the simple-pole
self-contraction persists. The critical value $k = -h^\vee$
gives $\kappa = 0$
(critical level, where this Shapovalov proof does not apply and
the Feigin--Frenkel centre changes the ambient problem). The shadow class is~L:
the cubic shadow $S_3 \neq 0$ detects the Lie bracket, but
$S_4 = 0$, so the tower terminates at degree~$3$.
The Koszul conductor is
$K = \kappa(V_k(\fg)) + \kappa(V_{k'}(\fg)) = 0$
where $k' = -2h^\vee - k$.

\subsection{Virasoro}

The Virasoro algebra $\Vir_c$ has OPE pole order~$4$ and is not
quadratic in any classical sense. Its Koszulness is proved
through shadow-tower formality~(iv). The PBW spectral sequence
for $\Vir_c$ has higher differentials that could a priori be
nonzero; the proof that they vanish uses the recursive structure
of the Maurer--Cartan element.

At genus zero, the shadow obstruction tower has
$o_{r+1}^{(0)} = 0$ for all $r \geq 2$: the transferred
$\Linfty$-brackets $\ell_{r+1}^{(0)}$ vanish because the
Virasoro vacuum module satisfies Kac determinant nonvanishing at
generic~$c$. Shadow-tower formality at genus zero is
equivalence~(iv)$\Leftrightarrow$(ii): the PBW spectral sequence
collapses.

The modular characteristic is $\kappa(\Vir_c) = c/2$, and the
shadow class is~M: $S_4 = 10/[c(5c+22)] \neq 0$ for generic~$c$,
so the all-genera tower is infinite.
The Koszul dual is $\Vir_c^! = \Vir_{26-c}$, giving
$K = c/2 + (26-c)/2 = 13$.
The self-dual point is $c = 13$ (not $c = 26$).

\subsection{Principal $\cW$-algebras}

The principal $\cW$-algebra $\cW_N$ at level~$k$ is obtained
from $V_k(\mathfrak{sl}_N)$ by Drinfeld--Sokolov reduction with
respect to the principal nilpotent $f_{\mathrm{prin}}$. The
Koszulness of $\cW_N$ is derived from the Koszulness of the
parent affine KM algebra through the functoriality of DS
reduction with respect to the bar complex.

DS reduction preserves the Shapovalov criterion: it sends the
Kac--Shapovalov form on $V_k(\mathfrak{sl}_N)$ to the
$\cW$-algebra Shapovalov form on $\cW_N$, and nonvanishing of
the former at generic level implies nonvanishing of the latter.
This gives Koszulness of $\cW_N$ through
equivalence~(xiii). At $N = 2$, this recovers the Virasoro
result through the DS path from $\mathfrak{sl}_2$.

The generators of $\cW_N$ have conformal weights
$\{2, 3, \ldots, N\}$: that is $N - 1$ generators, with highest
weight~$N$.
The modular characteristic is
$\kappa(\cW_N) = c \cdot (H_N - 1)$
where $H_N = \sum_{j=1}^{N} 1/j$. Sanity check at $N = 2$:
$H_2 - 1 = 3/2 - 1 = 1/2$, so $\kappa(\cW_2) = c/2$,
matching $\kappa(\Vir_c)$. At $N = 3$: $H_3 - 1 = 11/6 - 1
= 5/6$, so $\kappa(\cW_3) = 5c/6$.

\subsection{The $\beta\gamma$ system}

The $\beta\gamma$ system at conformal weight~$\lambda$ has central
charge $c_{\beta\gamma}(\lambda) = 2(6\lambda^2 - 6\lambda + 1)$.
Its bar complex has a contact-term tower that terminates at
degree~$4$ (shadow class~C), so the transferred $\Ainfty$-products
$m_n$ satisfy $m_n = 0$ for $n \geq 5$. But Koszulness requires
$m_n = 0$ for $n \geq 3$.

The vanishing of $m_3$ and $m_4$ is proved by explicit computation:
the genus-zero shadow obstructions $o_3^{(0)} = o_4^{(0)} = 0$ at
generic conformal weight. This establishes $\Ainfty$-formality,
equivalence~(iii). The $\beta\gamma$ system is the simplest
family in class~C, the class where the tower is finite but
nontrivial ($S_3 \neq 0$, $S_4 \neq 0$, $S_r = 0$ for
$r \geq 5$).

\subsection{Bershadsky--Polyakov}

The Bershadsky--Polyakov algebra
$\cW^{(2)}_k(\mathfrak{sl}_3)$ is the non-principal
$\cW$-algebra obtained by DS reduction of
$V_k(\mathfrak{sl}_3)$ with respect to the minimal nilpotent.
It has generators at conformal weights $\{1, 3/2, 3/2, 2\}$
and central charge, following Feigin--Kac--Roan--Wakimoto:
\[
 c(\mathrm{BP}_k) \;=\; 2 - \frac{24(k+1)^2}{k+3}.
\]
The formula is singular at the critical level $k = -3 = -h^\vee(\mathfrak{sl}_3)$, where the Sugawara denominator $k+h^\vee$ of the parent $V_k(\mathfrak{sl}_3)$ vanishes; the self-dual point $k=-3$ of the involution $k \mapsto -k - 6$ coincides with this critical locus, and the value $\kappa(\mathrm{BP}_{-3}) = 49/3$ quoted below is the \emph{principal value} of the symmetric limit across the pole, not a regular function value. The polynomial identity $K_{\mathrm{BP}}(k) := c(k) + c(-k-6) \equiv 196$ holds in $\mathbb{Q}(k)$, and $c(k) - 98$ is an odd function of $(k+3)$.

Koszulness is established through Ext diagonal
vanishing~(xii): the bigraded Ext groups
$\Ext^{p,q}(\omega, \omega)$ are concentrated on the diagonal
$p = q$ at generic level. The proof uses the non-principal DS
reduction from $V_k(\mathfrak{sl}_3)$: the DS functor sends
the Kac--Shapovalov form of the parent to the
Bershadsky--Polyakov Shapovalov form, and nonvanishing at
generic level gives diagonal concentration.

The modular characteristic is $\kappa(\cW^{(2)}_k) = c(k)/6$,
with anomaly ratio $\varrho = 1/6$. The Koszul complementarity
conductor is
$K_{\mathrm{BP}} = \kappa(\cW^{(2)}_k) + \kappa(\cW^{(2)}_{k'})
= 98/3$,
where $k' = -k - 6$ is the dual level.
The self-dual level is $k = -3$, at which
$\kappa = 49/3$.


% ================================================================
\section{The admissible-level question}
\label{sec:admissible}
% ================================================================

The standard landscape verification of
\S\ref{sec:landscape} treats the universal enveloping VOA
$V_k(\fg)$ at generic level. At admissible levels
$k = -h^\vee + p/q$ (coprime $p \geq h^\vee$, $q \geq 2$), the
simple quotient $L_k(\fg) = V_k(\fg)/I_k$ is a rational
vertex algebra. The question is whether the simple quotient
remains Koszul; the universal-algebra argument above does not
settle that quotient.

\subsection{Rank one: $\mathfrak{sl}_2$}

\begin{proposition}[$\mathfrak{sl}_2$ at admissible level]
\label{prop:kfc-sl2-admissible}
The simple quotient $L_k(\mathfrak{sl}_2)$ is chirally Koszul
at every admissible level
$k = -2 + p/q$ with $p \geq 2$, $q \geq 2$, $\gcd(p,q) = 1$.
\end{proposition}

\begin{proof}[Proof sketch]
The ideal $I_k \subset V_k(\mathfrak{sl}_2)$ is generated by a
single null vector at conformal weight $pq$. Since
$\dim \mathfrak{sl}_2 = 3$, the PBW spectral sequence for
$L_k(\mathfrak{sl}_2)$ has at most two nonzero columns at each
conformal weight: the Cartan direction $h$ and the root
directions $e, f$. The root generators are paired by
$[e,f] = h$ on the $E_1$ page, and the Cartan generator
survives as a single class in $H^1$.

The null vector at weight $pq$ introduces a relation in the bar
complex, but this relation is confined to a single PBW column: it
does not create off-diagonal classes in the $E_2$ page. The
PBW spectral sequence collapses at $\Etwo$ because the
dimensional constraint ($\dim \mathfrak{sl}_2 = 3$, hence at
most $2$ nonzero columns in the relevant range) leaves no room
for higher differentials $d_r$ with $r \geq 2$.
\end{proof}

\begin{remark}[Mechanism]
\label{rem:kfc-sl2-mechanism}
The proof for $\mathfrak{sl}_2$ uses the small-rank accident: three
generators produce at most two nonzero PBW columns at each weight,
and a differential $d_r$ for $r \geq 2$ needs at least $r + 1$
nonzero columns as domain and target. At rank one no
domain-target pair remains. This is specific to
$\mathfrak{sl}_2$; it fails at every higher rank.
\end{remark}

\subsection{Rank two and beyond: the obstruction}

\begin{conjecture}[Admissible Koszul rank obstruction]
\label{conj:kfc-admissible-rank}
For every simple Lie algebra $\fg$ of rank $r \geq 2$ and every
admissible level $k = p/q - h^\vee$ with denominator $q \geq 3$,
the simple quotient $L_k(\fg)$ has
$\dim H^2(\barBch(L_k(\fg))) = \mathrm{rank}(\fg)$.
Then $L_k(\fg)$ is not chirally Koszul.
\end{conjecture}

The mechanism is the abelian Cartan subalgebra
$\fh \subset \fg$. On the $E_1$ page of the PBW spectral
sequence for $L_k(\fg)$, the Poisson differential kills the root
sequence for $L_k(\fg)$, the Poisson differential pairs the root
generators through $[e_\alpha, f_\alpha] = h_\alpha$, but the
Cartan generators $h_1, \ldots, h_r$ survive because $\fh$ is
abelian. When $r = 1$, the surviving class is confined to a
single column and cannot support a higher differential. When
$r \geq 2$, the surviving $r$-dimensional space admits
nontrivial higher differentials sourced by the multi-weight
coupling between null vectors at different roots: the null
vectors interact through the Cartan matrix, and the
interactions generate off-diagonal bar cohomology.

\begin{remark}[Evidence]
\label{rem:kfc-admissible-evidence}
For $\fg = \mathfrak{sl}_3$ at admissible level with $q = 3$:
$\dim H^2(\barBch(L_k(\mathfrak{sl}_3))) = 2$, confirmed by
direct computation. The two classes correspond to the two simple
roots. For $\fg = \mathfrak{sl}_4$: predicted $H^2 = 3$.
The pattern is $\dim H^2 = \mathrm{rank}(\fg)$.
\end{remark}

\begin{remark}[Physical content]
\label{rem:kfc-admissible-physical}
Admissible-level simple quotients are the vertex algebras of
rational conformal field theories (Arakawa). If
Conjecture~\ref{conj:kfc-admissible-rank} holds, then
bar--cobar inversion fails for these quotients at rank
$\geq 2$: the null vectors that enforce rationality of the
representation category simultaneously obstruct PBW
concentration. Koszulness and rationality would stand in
tension for $\mathrm{rank} \geq 2$: the algebraic optimality
of the bar complex (Koszulness) is incompatible with the
representation-theoretic rigidity of the simple quotient
(rationality). At rank one, both coexist.
\end{remark}


% ================================================================
\section{The shadow algebra and BV/BRST}
\label{sec:shadow-bv}
% ================================================================

The Koszul property controls the bar complex at genus zero. At
higher genus, the bar differential squares to a nonzero
curvature $d^2_{\mathrm{fib}} = \kappa(\cA) \cdot \omega_g$, and
the Maurer--Cartan element $\Theta_\cA$ is the datum that
measures the obstruction to extending genus-zero
results across the modular tower. This section records two
structures that accompany the Koszul property: the shadow
algebra (the graded Lie algebra controlling the MC recursion)
and the BV/bar comparison (the genus-by-genus identification
with the BRST complex).

\subsection{The shadow algebra}

The shadow invariants $S_r(\cA)$ for $r \geq 2$ are the
degree-$r$ projections of $\Theta_\cA$. They do not form a ring:
the natural algebraic structure is a \emph{graded Lie bracket},
inherited from the convolution Lie algebra
$\fg^{\mathrm{mod}}_\cA$.

\begin{definition}[Shadow algebra]
\label{def:kfc-shadow-algebra}
The \emph{shadow algebra} of a chiral algebra~$\cA$ is the
bigraded Lie algebra
\[
\cA^{\mathrm{sh}} \;=\;
\bigoplus_{r \geq 2,\; g \geq 0} \cA^{\mathrm{sh}}_{r,g}
\]
where $\cA^{\mathrm{sh}}_{r,g}$ is the degree-$r$, genus-$g$
component of the modular convolution algebra
$\fg^{\mathrm{mod}}_\cA$. The Lie bracket
$[-,-]^{\mathrm{sh}} \colon
\cA^{\mathrm{sh}}_{r_1,g_1} \otimes \cA^{\mathrm{sh}}_{r_2,g_2}
\to \cA^{\mathrm{sh}}_{r_1+r_2-2, g_1+g_2}$
descends from the convolution bracket and drives the
Maurer--Cartan recursion: $S_{r+1}$ is determined by
$S_2, \ldots, S_r$ through
\[
d\,S_{r+1} + \tfrac{1}{2}
\sum_{\substack{r_1 + r_2 = r+2 \\ r_1, r_2 \geq 2}}
[S_{r_1}, S_{r_2}]^{\mathrm{sh}} = 0.
\]
\end{definition}

\begin{remark}[Lie-algebra structure]
\label{warn:kfc-shadow-not-ring}
The shadow algebra is a graded Lie algebra, \emph{not} a
graded-commutative ring. There is no shadow product
$S_r \cdot S_s$; there is a shadow bracket
$[S_r, S_s]^{\mathrm{sh}}$. The distinction is structural:
the Jacobi identity of $[-,-]^{\mathrm{sh}}$ controls the
solvability of the MC recursion, and the depth classification
G/L/C/M records where this recursion first produces nonzero
data.
\end{remark}

The G/L/C/M classification follows from the shadow algebra:
\begin{itemize}
\item Class~G ($r_{\max} = 2$): $S_r = 0$ for $r \geq 3$.
  The shadow algebra is abelian. Heisenberg, lattice algebras.
\item Class~L ($r_{\max} = 3$): $S_3 \neq 0$, $S_r = 0$ for
  $r \geq 4$. The cubic shadow detects the Lie bracket. Affine
  Kac--Moody at non-critical level.
\item Class~C ($r_{\max} = 4$): $S_3 \neq 0$, $S_4 \neq 0$,
  $S_r = 0$ for $r \geq 5$. Contact-term termination. The
  $\beta\gamma$ system.
\item Class~M ($r_{\max} = \infty$): $S_r \neq 0$ for
  infinitely many~$r$. Virasoro, $\cW$-algebras.
\end{itemize}
Every class is chirally Koszul. The depth classification is
orthogonal to Koszulness: it measures complexity within the
Koszul world, not failure of Koszulness.

\subsection{BV/bar identification in the coderived category}

The BV/BRST complex $C^\bullet_{\mathrm{BV}}(\cA, \Sigma_g)$
on a genus-$g$ surface $\Sigma_g$ computes the same factorization
homology as the bar complex $B^{(g)}(\cA)$, but through a
different chain model: the BV propagator is the Green's function
on $\Sigma_g$, while the bar propagator is the logarithmic form
$d\log E(z,w)$ (weight~$1$, the prime form logarithmic
derivative).

\begin{theorem}[BV$=$bar in the coderived category]
\label{thm:kfc-bv-bar}
Let $\cA$ be a chirally Koszul algebra.
\begin{enumerate}[label=\textup{(\roman*)}]
\item For every genus $g \geq 0$,
  $C^\bullet_{\mathrm{BV}}(\cA, \Sigma_g)
  \simeq_{D^{\mathrm{co}}} B^{(g)}(\cA)$
  in the coderived category of curved factorization models.
\item If $\cA$ is class~G or class~L, the comparison morphism
  $f_g$ is a chain-level weak equivalence.
\item If $\cA$ is class~C and harmonic decoupling holds,
  $f_g$ is a chain-level weak equivalence.
\item If $\cA$ is class~M, the comparison cone is coacyclic:
  $f_g$ is a coderived equivalence but \emph{not} a chain-level
  quasi-isomorphism. The harmonic discrepancy at degree~$r$ is
  $\delta_r^{\mathrm{harm}} = c_r(\cA) \cdot
  m_0^{\lfloor r/2 \rfloor - 1}$,
  a sum of curvature powers whose cone is exact in the coderived
  sense.
\end{enumerate}
\end{theorem}

\begin{proof}[Proof sketch]
The comparison morphism $f_g$ is constructed by replacing the BV
propagator with its Hodge decomposition relative to the bar
propagator. At genus~$0$, the Hodge decomposition is trivial (no
harmonic forms on $\bP^1$), and $f_0$ is a chain-level
quasi-isomorphism for all classes.

At genus $g \geq 1$, the Hodge decomposition produces a harmonic
correction proportional to the curvature
$m_0 = \kappa(\cA) \cdot \omega_g$. For class~G,
$\kappa = k$ and $S_r = 0$ for $r \geq 3$: the harmonic
correction at each degree involves only the curvature itself,
and the comparison is exact. For class~L, the cubic shadow
contributes a single additional correction at degree~$3$,
which vanishes by the Jacobi identity of the shadow bracket.

For class~M, the harmonic corrections are nonzero at
every degree $r \geq 4$. Each correction
$\delta_r^{\mathrm{harm}}$ is a polynomial in the curvature
$m_0$; the sum defines the comparison cone. This cone is
\emph{not} acyclic in the ordinary sense (the harmonic
terms persist at each degree), but it \emph{is} coacyclic:
it belongs to the acyclic class of the coderived category.
The proof uses the characterization of coacyclic objects as
those whose totalizations are contractible in the complete
filtered sense~\cite[\S3.7]{PP}.
\end{proof}

\begin{remark}[Class~M: chain-level false]
\label{rem:kfc-class-m-chain-false}
The failure of the chain-level BV/bar comparison for class~M
is genuine: it is not an artifact of a suboptimal chain model.
The Virasoro algebra has
$\delta_4^{\mathrm{harm}} = c_4(\Vir_c) \cdot m_0 \neq 0$
for $c \neq 0$, and $c_4$ is determined by the shadow
invariant $S_4 = 10/[c(5c+22)]$. Since $S_4 \neq 0$ for
generic~$c$, the harmonic discrepancy is nonzero, and no
chain-level comparison exists. The coderived category is
the correct ambient category for BV/bar at class~M.
\end{remark}

\begin{remark}[Scope summary]
\label{rem:kfc-bv-scope}
BV$=$bar in the coderived category is proved unconditionally for
all chirally Koszul algebras. The chain-level strengthening holds
for class~G and class~L (unconditionally), for class~C
(conditionally on harmonic decoupling), and fails for class~M.
The coderived formulation is the natural one: it accommodates
the curvature $d^2_{\mathrm{fib}} = \kappa \cdot \omega_g$
inside the coderived acyclic class.
\end{remark}


% ================================================================
\section{Structural summary}
\label{sec:first-principles-triples}
% ================================================================

The classical KZ criteria detect chiral Koszulness on the core PBW
component. The supplementary criteria are chart functors attached
to other associator base points. Their overlap is controlled by
Tamarkin transfer on finite PBW-Artin components carrying the same
Koszul witness.

For Virasoro at generic central charge, the Hodge--Betti chart
$U_{12}$ proves Koszulness using the Feigin--Fuchs BGG resolution
and the Kac determinant. Bar--cobar inversion is not an input to
that proof.

For Yangians, the associative Polishchuk--Positselski PBW theorem
supplies an associative Koszul chart. On the comparison component
there is an open embedding
$\cM_{\mathrm{Kosz}}^{\mathrm{assoc}}(Y_\hbar) \hookrightarrow
 \cM_{\mathrm{Kosz}}^{\mathrm{chiral}}(Y_\hbar)$
(Theorem~\ref{thm:yangian-chart-inclusion}). Chart-surjectivity
is an isomorphism at $\mathfrak{sl}_2$ and remains open in higher
rank.

% ================================================================
\bigskip

\noindent
The Koszul locus is the locus on which the bar construction is
invertible, the spectral sequence collapses, the resolution is
linear, and the scalar projection of the degree-two
Maurer--Cartan element is controlled by $\kappa$. Fourteen chart
functors detect this locus on their stated components.
The standard generic rows lie on the Koszul locus.
Shadow depth (the G/L/C/M classification) is an orthogonal invariant:
it classifies complexity within the Koszul world, not proximity to
the Koszul boundary. The admissible-level question (whether simple
quotients at rank $\geq 2$ remain Koszul) and the $\cD$-module
purity converse (whether Koszulness implies Hodge-theoretic purity
beyond affine Kac--Moody) are the two open frontiers where the
boundary of the Koszul locus has not yet been mapped.


% ================================================================
% Bibliography
% ================================================================

\begin{thebibliography}{99}

\bibitem{AT12}
A.~Alekseev and C.~Torossian.
\newblock The Kashiwara--Vergne conjecture and Drinfeld's
 associators.
\newblock \emph{Ann.~of Math.} \textbf{175} (2012), 415--463.

\bibitem{BD04}
A.~Beilinson and V.~Drinfeld.
\newblock \emph{Chiral Algebras}.
\newblock AMS Colloquium Publications, vol.~51, 2004.

\bibitem{Deligne89}
P.~Deligne.
\newblock Le groupe fondamental de la droite projective moins
 trois points.
\newblock In \emph{Galois Groups over~$\bQ$}, MSRI Publ.~16,
 Springer, 1989, pp.~79--297.

\bibitem{Drinfeld90}
V.~G.~Drinfeld.
\newblock On quasitriangular quasi-Hopf algebras and a group
 closely connected with $\mathrm{Gal}(\overline{\bQ}/\bQ)$.
\newblock \emph{Leningrad Math.~J.} \textbf{2} (1991), 829--860.

\bibitem{Enriquez14}
B.~Enriquez.
\newblock Elliptic associators.
\newblock \emph{Selecta Math.} \textbf{20} (2014), 491--584.

\bibitem{FeiginFuchs90}
B.~L.~Feigin and D.~B.~Fuchs.
\newblock Representations of the Virasoro algebra.
\newblock In \emph{Representation of Lie Groups and Related
 Topics}, Adv.~Stud.~Contemp.~Math.~\textbf{7}, Gordon and
 Breach, 1990, pp.~465--554.

\bibitem{Furusho10}
H.~Furusho.
\newblock Pentagon and hexagon equations.
\newblock \emph{Ann.~of Math.} \textbf{171} (2010), 545--556.

\bibitem{KacBook}
V.~G.~Kac.
\newblock \emph{Infinite Dimensional Lie Algebras}.
\newblock Third edition, Cambridge University Press, 1990.

\bibitem{Kontsevich93}
M.~Kontsevich.
\newblock Vassiliev's knot invariants.
\newblock In \emph{I.~M.~Gel'fand Seminar}, Adv.~Soviet
 Math.~\textbf{16}, AMS, 1993, pp.~137--150.

\bibitem{Kontsevich99}
M.~Kontsevich.
\newblock Operads and motives in deformation quantization.
\newblock \emph{Lett.~Math.~Phys.} \textbf{48} (1999), 35--72.

\bibitem{RCW82}
A.~Rocha-Caridi and N.~R.~Wallach.
\newblock Characters of irreducible representations of the
 Virasoro algebra.
\newblock \emph{Math.~Z.} \textbf{185} (1984), 1--21.

\bibitem{Tamarkin03}
D.~Tamarkin.
\newblock Formality of chain operad of little discs.
\newblock \emph{Lett.~Math.~Phys.} \textbf{66} (2003), 65--72.

\bibitem{BGS}
A.~Beilinson, V.~Ginzburg, and W.~Soergel.
\newblock Koszul duality patterns in representation theory.
\newblock \emph{J.~Amer.~Math.~Soc.} \textbf{9} (1996), 473--527.

\bibitem{CG}
K.~Costello and O.~Gwilliam.
\newblock \emph{Factorization algebras in quantum field theory},
 vols.~I--II.
\newblock Cambridge University Press, 2017/2021.

\bibitem{FG12}
J.~Francis and D.~Gaitsgory.
\newblock Chiral Koszul duality.
\newblock \emph{Selecta Math.~(N.S.)} \textbf{18} (2012), 27--87.

\bibitem{FultonMacPherson94}
W.~Fulton and R.~MacPherson.
\newblock A compactification of configuration spaces.
\newblock \emph{Ann.~of Math.} \textbf{139} (1994), 183--225.

\bibitem{GL76}
H.~Garland and J.~Lepowsky.
\newblock Lie algebra homology and the Macdonald--Kac formulas.
\newblock \emph{Invent.~Math.} \textbf{34} (1976), 37--76.

\bibitem{GR17}
D.~Gaitsgory and N.~Rozenblyum.
\newblock \emph{A study in derived algebraic geometry}.
\newblock American Mathematical Society, 2017.

\bibitem{HA}
J.~Lurie.
\newblock \emph{Higher Algebra}.
\newblock 2017. Available at
 \url{https://www.math.ias.edu/~lurie/}.

\bibitem{KK79}
V.~Kac and D.~Kazhdan.
\newblock Structure of representations with highest weight of
 infinite-dimensional Lie algebras.
\newblock \emph{Adv.~in Math.} \textbf{34} (1979), 97--108.

\bibitem{Keller}
B.~Keller.
\newblock $A$-infinity algebras, modules and functor categories.
\newblock In \emph{Trends in Representation Theory of Algebras
 and Related Topics}, Contemp.~Math.~\textbf{406} (2006), 67--93.

\bibitem{LV12}
J.-L.~Loday and B.~Vallette.
\newblock \emph{Algebraic Operads}.
\newblock Grundlehren der math.~Wissenschaften, vol.~346,
 Springer, 2012.

\bibitem{LorgatMKD1}
R.~Lorgat.
\newblock \emph{Modular Koszul Duality}.
\newblock In preparation, 2026.

\bibitem{PP}
A.~Polishchuk and L.~Positselski.
\newblock \emph{Quadratic Algebras}.
\newblock University Lecture Series, vol.~37, AMS, 2005.

\bibitem{Priddy}
S.~Priddy.
\newblock Koszul resolutions.
\newblock \emph{Trans.~Amer.~Math.~Soc.} \textbf{152} (1970), 39--60.

\bibitem{Saito}
M.~Saito.
\newblock Modules de Hodge polarisables.
\newblock \emph{Publ.~Res.~Inst.~Math.~Sci.} \textbf{24} (1988), 849--995.

\bibitem{STS83}
M.~Semenov-Tian-Shansky.
\newblock What is a classical $r$-matrix?
\newblock \emph{Funct.~Anal.~Appl.} \textbf{17} (1983), 259--272.

\end{thebibliography}

\end{document}
